% ============================================================
\section{Analysis Suite: Carbon, Production, Lifecycle, Distribution, Reliability, and Time Series}
\label{sec:analysis}
% ============================================================

This chapter documents the analytical capability modules of the
\texttt{hacdcpf} library.  Following a module restructuring, each family of
analytical algorithms now lives in its own peer directory rather than a shared
\texttt{analysis/} facade.  The current layout is:

\begin{center}
\begin{tabular}{ll}
\hline\hline
\textbf{Module directory} & \textbf{Contents} \\
\hline
\texttt{include/hacdcpf/carbon\_analysis} / \texttt{src/carbon\_analysis} & Carbon emissions analysis (\texttt{compute\_carbon\_analysis}) \\
\texttt{include/hacdcpf/time\_series} / \texttt{src/time\_series} & UC/OPF/PF multi-period pipeline, annual production simulation, lifecycle simulation \\
\texttt{include/hacdcpf/reliability} / \texttt{src/reliability} & Monte Carlo reliability assessment \\
\texttt{include/hacdcpf/resilience} / \texttt{src/resilience} & Distribution resilience assessment \\
\texttt{include/hacdcpf/power\_flow} / \texttt{src/power\_flow} & Core Newton PF solvers plus the three-phase distribution power-flow NR helpers \\
\hline\hline
\end{tabular}
\end{center}

All modules are declared in namespace \texttt{hacdcpf::analysis} (or the
top-level \texttt{hacdcpf} namespace for shared data types) and exposed
through \texttt{include/hacdcpf/api/hacdcpf.hpp}.

% ------------------------------------------------------------
\subsection{Carbon Emissions Analysis}
\label{sec:carbon}
% ------------------------------------------------------------

Declared in \texttt{carbon\_analysis/carbon\_analysis.hpp},
implemented in \texttt{src/carbon\_analysis/carbon\_analysis.cpp}.

The primary entry point is \texttt{compute\_carbon\_analysis}, which has two
overloads:
\begin{verbatim}
// Overload 1: supply an already-computed power-flow result.
CarbonAnalysisResult compute_carbon_analysis(
    const HybridPowerSystem& sys,
    const PowerFlowResult&   pf_result,
    const CarbonAnalysisOptions& opt = {});

// Overload 2: run power flow internally, then analyse.
CarbonAnalysisResult compute_carbon_analysis(
    const HybridPowerSystem& sys,
    const PowerFlowOptions&  pf_opt  = {},
    const CarbonAnalysisOptions& ca_opt = {});
\end{verbatim}
A backward-compatible inline wrapper \texttt{run\_carbon\_analysis(sys, opt)}
calls overload~2 with default \texttt{PowerFlowOptions}.

\texttt{compute\_carbon\_analysis} performs power-flow-based emissions
attribution via two independent methods:
\begin{enumerate}
  \item \textbf{Proportional tracing} --- Kirchhoff-based upstream tracing that
        allocates each generator's output to loads and losses in proportion to
        power flow direction.
  \item \textbf{Nodal matrix method} --- Builds and solves a linear system
        $\mathbf{A}\,\mathbf{c} = \mathbf{q}$ where $c_i$ is the carbon
        intensity at bus $i$ (tCO\textsubscript{2}/MWh).
\end{enumerate}

\subsubsection{Mathematical Formulation}

\paragraph{Proportional Tracing.}
Let $\mathcal{G}_i$ denote generators injecting into bus $i$, and let
$\mathcal{U}(i) = \{j : P_{j\to i} > 0\}$ be the set of upstream buses
with active inflow.  The total inflow at bus $i$ is:
\begin{equation}
  T_i = \sum_{g \in \mathcal{G}_i} P_{g,i} + \sum_{j \in \mathcal{U}(i)} P_{j\to i}.
  \label{eq:carbon-inflow}
\end{equation}
The carbon intensity at bus $i$ is computed by the upstream tracing rule:
\begin{equation}
  c_i = \frac{\sum_{g \in \mathcal{G}_i} e_g\,P_{g,i}
              + \sum_{j \in \mathcal{U}(i)} c_j\,P_{j\to i}}{T_i},
  \label{eq:carbon-tracing}
\end{equation}
where $e_g$ (tCO\textsubscript{2}/MWh) is the emission factor of generator $g$.
The system is solved by topological ordering from generator buses outward.
Mass-balance verification requires:
\begin{equation}
  \left|\sum_{g} e_g\,P_g - \sum_{i\in\mathcal{L}} c_i\,P_i^L
        - \sum_{(i,j)} c_i\,P_{ij}^{loss}\right| \leq \varepsilon_{bal}.
\end{equation}

\paragraph{Nodal Matrix Method.}
Equations~\eqref{eq:carbon-inflow}--\eqref{eq:carbon-tracing} can be
recast as a linear system $\mathbf{A}\mathbf{c} = \mathbf{q}$ with:
\begin{align}
  A_{ii} &= T_i = \sum_{g\in\mathcal{G}_i} P_{g,i}
                  + \sum_{j\in\mathcal{U}(i)} P_{j\to i}, \\
  A_{ij} &= -P_{j\to i} \quad (j \neq i,\; j \in \mathcal{U}(i)), \\
  q_i    &= \sum_{g\in\mathcal{G}_i} e_g\,P_{g,i}.
\end{align}
The matrix $\mathbf{A}$ is diagonally dominant and its solution $\mathbf{c}$
is identical to the tracing result when $\mathbf{A}$ is non-singular.
The residual $\|\mathbf{A}\mathbf{c} - \mathbf{q}\|_\infty$ is stored in
\texttt{CarbonAnalysisResult::matrix\_residual}.

\subsubsection{Options and Result Types}

\begin{center}
\begin{tabular}{lll}
\hline\hline
\textbf{Field (CarbonAnalysisOptions)} & \textbf{Default} & \textbf{Description} \\
\hline
\texttt{min\_contribution\_mw} & \texttt{1e-6} & Minimum generator contribution to track (MW) \\
\texttt{verification\_tol} & \texttt{1e-4} & Relative balance tolerance for tracing verification \\
\texttt{max\_tracing\_depth} & 0 & BFS depth limit (0~=~unlimited) \\
\texttt{loss\_allocation\_alpha} & 0.5 & Loss split: $\alpha$ at sending end, $(1-\alpha)$ at receiving end \\
\texttt{regularization\_eps} & \texttt{1e-8} & Diagonal regulariser for near-zero pivot buses \\
\texttt{verbose} & false & Emit debug-level logging during analysis \\
\hline\hline
\end{tabular}
\end{center}

\texttt{CarbonAnalysisResult} aggregates per-component results:

\begin{longtable}{ll}
\hline\hline
\textbf{Field} & \textbf{Description} \\
\hline\endhead\hline\endfoot
\texttt{tracing\_verified} & Tracing mass-balance check passed \\
\texttt{matrix\_solved} & Matrix method converged \\
\texttt{matrix\_residual} & Residual of the linear system \\
\texttt{load\_carbon} & Per-AC-load: demand, carbon intensity, total emissions, supply mix \\
\texttt{dc\_load\_carbon} & Same for DC loads \\
\texttt{branch\_carbon} & Per-AC-branch: loss, carbon intensity, generator attribution \\
\texttt{dc\_branch\_carbon} & Same for DC branches \\
\texttt{bus\_carbon} & Per-AC-bus carbon intensity (tCO\textsubscript{2}/MWh) \\
\texttt{dc\_bus\_carbon} & Same for DC buses \\
\texttt{vsc\_carbon} & Per-VSC: input/output power, losses, carbon attribution \\
\texttt{dcdc\_carbon} & Per-DC-DC converter: same \\
\texttt{storage\_carbon} & Per-storage: SoC, stored energy, emissions \\
\texttt{energy\_router\_carbon} & Per-router: port-level power and emissions \\
\texttt{generator\_loss\_allocation} & Per-generator share of network losses \\
\texttt{tracing\_summary} & \texttt{EmissionsSummary}: generation, load, loss, balance error \\
\texttt{matrix\_summary} & Same from matrix method \\
\end{longtable}

\subsubsection{EmissionsSummary}
\begin{center}
\begin{tabular}{ll}
\hline\hline
\textbf{Field} & \textbf{Description} \\
\hline
\texttt{total\_generation\_emissions\_tco2} & Weighted sum of generator emission rates \\
\texttt{total\_load\_emissions\_tco2} & Emissions attributed to all loads \\
\texttt{total\_loss\_emissions\_tco2} & Emissions attributed to network losses \\
\texttt{balance\_error\_tco2} & $|E_{gen} - E_{load} - E_{loss}|$ \\
\texttt{balance\_error\_pct} & Relative balance error (\%) \\
\hline\hline
\end{tabular}
\end{center}

\subsubsection{Loss Allocation and Directed Flow Edges}
\label{sec:carbon-flow-edges}

Every branch in the hybrid AC/DC network is represented as a directed
\emph{flow edge} $(i \to j)$ carrying three power quantities:
\begin{equation}
  P_{i\to j}^{\text{send}},\quad
  P_{i\to j}^{\text{recv}},\quad
  \ell_{ij} = P_{i\to j}^{\text{send}} - P_{i\to j}^{\text{recv}} \geq 0.
\end{equation}
For AC branches the direction follows the sign of $P_{ij}^{\text{from}}$;
for DC branches $P^{\text{from}} = V_i(V_i - V_j)/R$.  VSC converters,
DC/DC converters and Energy Routers contribute edges whose
$(P^{\text{send}}, P^{\text{recv}})$ are taken directly from the
power-flow transfer objects.

The loss-allocation parameter $\alpha \in [0,1]$ (default $0.5$) splits
each branch loss between the sending and receiving nodes.  The power
\emph{transported} from $i$ to $j$ (i.e.\ available to downstream buses)
is:
\begin{equation}
  P_{ij}^{\text{trans}} = \max\!\bigl(P_{i\to j}^{\text{send}}
                           - \alpha\,\ell_{ij},\;0\bigr).
  \label{eq:carbon-transported}
\end{equation}
For the nodal matrix method, the receiving-end power balance adds
$P_{ij}^{\text{trans}}$ to bus~$j$'s total inflow while attributing
$(1-\alpha)\ell_{ij}$ to bus~$j$'s outflow (as a local loss sink):
\begin{equation}
  P_{\text{in}}(j) \mathrel{+}= P_{ij}^{\text{trans}}, \qquad
  P_{\text{out}}(i) \mathrel{+}= P_{i\to j}^{\text{send}}, \qquad
  P_{\text{out}}(j) \mathrel{+}= (1-\alpha)\,\ell_{ij}.
\end{equation}
The matrix system then becomes:
\begin{equation}
  A_{ii} = P_{\text{out}}(i), \qquad
  A_{ij} = -P_{ji}^{\text{trans}} \quad (j \neq i),
  \qquad b_i = \sum_{g \in \mathcal{G}_i} e_g\,P_{g,i},
  \label{eq:carbon-matrix-full}
\end{equation}
where $P_{ji}^{\text{trans}}$ is the transported power on any edge
arriving at bus~$i$ from bus~$j$.  Isolated transit buses (zero load
and generation, $P_{\text{out}}(i) < \varepsilon_{\text{reg}}$) are
regularised by a Laplacian-average constraint using neighbouring bus
intensities.

\subsubsection{Storage State-of-Charge Carbon Intensity}
\label{sec:carbon-storage}

Battery storage is modelled as a carbon-carrying reservoir.  When a
storage unit charges ($p^{\text{ch}} > 0$, net inflow), the stored
energy carbon intensity is updated by energy-weighted mixing:
\begin{equation}
  \lambda_{t+1}^{\text{soc}} =
    \frac{E_t\,\lambda_t^{\text{soc}} + \Delta E_{\text{in}}\,c_b}{E_{t+1}},
  \label{eq:storage-soc-carbon}
\end{equation}
where $E_t = \text{SoC}_t \times E_{\text{rated}}$ is stored energy at
the start of the step, $\Delta E_{\text{in}} = \max(E_{t+1}-E_t, 0)$
is the incremental charge, and $c_b$ is the nodal carbon intensity of
the charging bus (from the proportional tracing or matrix method).
When the unit discharges, $\lambda^{\text{soc}}$ carries forward
unchanged and is used to attribute emissions to the discharged energy.

% ------------------------------------------------------------
\subsection{Annual Carbon Analysis}
\label{sec:annual-carbon}
% ------------------------------------------------------------

Declared in \texttt{annual\_carbon\_analysis.hpp},
implemented in \texttt{src/carbon\_analysis/annual\_carbon\_analysis.cpp}.
The primary entry point is:
\begin{verbatim}
AnnualCarbonAnalysisResult compute_annual_carbon_analysis(
    const HybridPowerSystem& sys,                   // or vector<HybridPowerSystem>
    const std::vector<PowerFlowResult>& pf_results,
    double step_duration_hr,
    const AnnualCarbonAnalysisOptions& options = {});
\end{verbatim}

\subsubsection{Multi-Step Aggregation}

Given $N_T$ converged time steps each of duration $\Delta t$ (hours), the
per-step emissions and energy at bus~$i$ are:
\begin{equation}
  E_i(t) = P_i^L(t)\,\Delta t, \qquad
  \epsilon_i(t) = c_i(t)\,E_i(t),
\end{equation}
where $c_i(t)$ is the nodal carbon intensity (tCO\textsubscript{2}/MWh)
from \texttt{compute\_carbon\_analysis} at step $t$.  Annual totals:
\begin{equation}
  E_i^{\text{ann}} = \sum_{t=1}^{N_T} E_i(t), \qquad
  \epsilon_i^{\text{ann}} = \sum_{t=1}^{N_T} \epsilon_i(t).
\end{equation}
Two distinct carbon-intensity averages are reported for each bus:
\begin{align}
  \bar{c}_i^{\text{load}} &= \frac{\epsilon_i^{\text{ann}}}{E_i^{\text{ann}}}
    = \frac{\sum_t c_i(t)\,E_i(t)}{\sum_t E_i(t)},
    \label{eq:annual-load-weighted} \\[4pt]
  \bar{c}_i^{\text{time}} &= \frac{\sum_t c_i(t)\,\Delta t}{N_T\,\Delta t}
    = \frac{1}{N_T}\sum_t c_i(t).
    \label{eq:annual-time-weighted}
\end{align}
$\bar{c}_i^{\text{load}}$ weights by energy consumed (equivalent to
average emission rate per kWh delivered) while $\bar{c}_i^{\text{time}}$
weights uniformly by time (average grid carbon intensity experienced over
the study period).

\subsubsection{Storage Carbon State Propagation}

When the multi-system overload is used (one \texttt{HybridPowerSystem}
per step, enabling topology and dispatch changes), the storage SoC
carbon intensity $\lambda^{\text{soc}}$ is carried forward between
consecutive systems according to Equation~\eqref{eq:storage-soc-carbon}.
This preserves the carbon provenance of energy stored in earlier periods.

\subsubsection{Green Energy Certificate (GEC) Accounting}
\label{sec:gec}

A user $u$ owns a portfolio of loads $\mathcal{L}_u$ and holds an annual
GEC quota $G_u$ (MWh) with a certificate carbon intensity $\lambda_G^u$
(tCO\textsubscript{2}/MWh) representing the renewable source.
\texttt{compute\_annual\_user\_gec\_accounting} distributes $G_u$ across
time steps pro-rata to the user's consumption share and computes net
emissions after GEC offsetting.

\paragraph{Step-level GEC allocation.}
The planned GEC fraction for step $t$ is proportional to the step's
energy share:
\begin{equation}
  \phi_{u,t} = \frac{E_{u,t}}{E_u^{\text{ann}}}, \qquad
  G_{u,t}^{\text{plan}} = G_u\,\phi_{u,t}.
\end{equation}
To ensure GECs do not exceed actual consumption, the allocated volume is
capped:
\begin{equation}
  G_{u,t} = \min\!\bigl(G_{u,t}^{\text{plan}},\; E_{u,t}\bigr).
  \label{eq:gec-allocated}
\end{equation}
Net emissions for step $t$ are then:
\begin{equation}
  \epsilon_{u,t}^{\text{net}} =
    \underbrace{(E_{u,t} - G_{u,t})\,c_u(t)}_{\text{grid-sourced emissions}}
    + \underbrace{G_{u,t}\,\lambda_G^u}_{\text{GEC emissions}},
  \label{eq:gec-net}
\end{equation}
where $c_u(t) = \epsilon_{u,t}^{\text{gross}} / E_{u,t}$ is the gross
nodal carbon intensity averaged over the user's loads.

\paragraph{Annual GEC summary.}
Annual totals follow directly:
\begin{equation}
  G_u^{\text{used}} = \sum_t G_{u,t}, \qquad
  G_u^{\text{unused}} = G_u - G_u^{\text{used}}.
\end{equation}
The avoided emissions and net intensity are:
\begin{align}
  \epsilon_u^{\text{avoided}} &= \epsilon_u^{\text{gross}} - \epsilon_u^{\text{net}}, \\
  \bar{c}_u^{\text{gross}}    &= \frac{\epsilon_u^{\text{gross}}}{E_u^{\text{ann}}}, \\
  \bar{c}_u^{\text{net}}      &= \frac{\epsilon_u^{\text{net}}}{E_u^{\text{ann}}}, \\
  \rho_u                      &= \frac{G_u^{\text{used}}}{E_u^{\text{ann}}},
  \label{eq:gec-coverage}
\end{align}
where $\rho_u \in [0,1]$ is the GEC coverage ratio (fraction of annual
consumption offset by certificates).

\begin{center}
\begin{tabular}{lll}
\hline\hline
\textbf{Field (AnnualCarbonAnalysisOptions)} & \textbf{Default} & \textbf{Description} \\
\hline
\texttt{carbon\_options} & (default) & \texttt{CarbonAnalysisOptions} forwarded to each step \\
\texttt{keep\_hourly\_bus\_intensity} & false & Retain per-step bus intensity matrix \\
\texttt{keep\_hourly\_load\_emissions} & false & Retain per-step load emissions matrix \\
\texttt{keep\_hourly\_load\_energy} & false & Retain per-step load energy matrix \\
\hline\hline
\end{tabular}
\end{center}

% ------------------------------------------------------------
\subsection{Distribution Network Power Flow}
\label{sec:dist-pf}
% ------------------------------------------------------------

Declared in \texttt{power\_flow/distribution\_power\_flow.hpp},
implemented in \texttt{src/power\_flow/three\_phase\_nr.cpp}.
The public distribution power-flow surface is intentionally three-phase only:
legacy single-phase radial/backward-forward sweep entry points and the old
\texttt{DistributionPFOptions}/\texttt{DistributionPFResult} contract are not
exported.  Distribution feeders are solved through the full abc-domain
Newton--Raphson interface:
\begin{verbatim}
ThreePhaseDPFResult solve_three_phase_nr(
    const ThreePhaseACSystem& sys,
    const ThreePhaseNROptions& opt = {});

ThreePhaseDPFResult solve_three_phase_nr(
    const HybridPowerSystem& sys,
    const ThreePhaseNROptions& opt = {});
\end{verbatim}
The same header also exports \texttt{build\_full\_ybus\_phase\_entries} for
phase-domain admittance inspection and regression tests.

\subsubsection{Three-Phase Newton-Raphson}

For an $n$-bus system with per-phase admittances, define the $3n$-element
state vector $\mathbf{x} = [\boldsymbol{\theta}^a, \boldsymbol{\theta}^b,
\boldsymbol{\theta}^c, \mathbf{V}_m^a, \mathbf{V}_m^b, \mathbf{V}_m^c]^\top$.
The three-phase nodal admittance matrix is:
\begin{equation}
  \mathbf{Y}^{abc}_{ij} = \begin{bmatrix}
    Y_{ij}^{aa} & Y_{ij}^{ab} & Y_{ij}^{ac} \\
    Y_{ij}^{ba} & Y_{ij}^{bb} & Y_{ij}^{bc} \\
    Y_{ij}^{ca} & Y_{ij}^{cb} & Y_{ij}^{cc}
  \end{bmatrix}.
\end{equation}
The power-mismatch residual at each phase-bus is:
\begin{equation}
  \mathbf{f}(\mathbf{x}) = \mathbf{S}^{inj} - \text{diag}(\mathbf{V}^{abc})
    \left(\mathbf{Y}^{abc}_{bus}\,\mathbf{V}^{abc}\right)^* = \mathbf{0},
\end{equation}
and the Newton step solves:
\begin{equation}
  \mathbf{J}^{(k)}\,\Delta\mathbf{x}^{(k)} = -\mathbf{f}\!\left(\mathbf{x}^{(k)}\right),
  \qquad \mathbf{x}^{(k+1)} = \mathbf{x}^{(k)} + \Delta\mathbf{x}^{(k)},
\end{equation}
where $\mathbf{J}^{(k)} = \partial\mathbf{f}/\partial\mathbf{x}|_{\mathbf{x}^{(k)}}$
is the $6n\times 6n$ Jacobian.  The voltage unbalance factor at bus $i$ is:
\begin{equation}
  \text{VUF}_i = \frac{|V_i^{neg}|}{|V_i^{pos}|}\times 100\,\%,
\end{equation}
where $V_i^{neg}$/$V_i^{pos}$ are the negative and positive sequence
components from Fortescue's symmetrical-components transformation.

\subsubsection{Three-Phase NR Options and Results}

\begin{center}
\begin{tabular}{llll}
\hline\hline
\textbf{Field (ThreePhaseNROptions)} & \textbf{Default} & \textbf{Units} & \textbf{Description} \\
\hline
\texttt{max\_iter} & 50 & --- & Maximum Newton iterations per control step \\
\texttt{max\_control\_iter} & 100 & --- & Maximum regulator/control iterations \\
\texttt{tol} & \(10^{-6}\) & pu & Residual convergence tolerance \\
\texttt{verbose} & false & --- & Emit iteration diagnostics \\
\texttt{include\_shunts} & true & --- & Include shunt admittances in the phase-domain Y-bus \\
\hline\hline
\end{tabular}
\end{center}

\begin{center}
\begin{tabular}{ll}
\hline\hline
\textbf{Field (ThreePhaseDPFResult)} & \textbf{Description} \\
\hline
\texttt{bus\_voltages} & Per-bus, per-phase voltage magnitudes and angles \\
\texttt{branch\_powers} & Per-line, per-phase sending-end active/reactive powers \\
\texttt{transformer\_terminal\_observations} & Per-terminal transformer voltage/current observations \\
\texttt{regulator\_states}, \texttt{regulator\_trace} & Final regulator state and control trace \\
\texttt{p/q\_loss\_*}, \texttt{total\_p/q\_loss\_*} & Per-phase and total losses \\
\texttt{max\_vuf\_percent} & Worst voltage-unbalance factor \\
\texttt{control\_converged, converged, iterations, residual} & Control and Newton convergence indicators \\
\hline\hline
\end{tabular}
\end{center}

% ------------------------------------------------------------
\subsection{Three-Phase Unbalanced Power Flow}
\label{sec:three-phase}
% ------------------------------------------------------------

Declared in \texttt{power\_flow/three\_phase.hpp}.  The primary solver
(\texttt{solve\_three\_phase(sys, opt)}) delegates to
\texttt{analysis::solve\_three\_phase\_nr()} and reports:

\begin{center}
\begin{tabular}{ll}
\hline\hline
\textbf{Field (ThreePhaseFlowResult)} & \textbf{Description} \\
\hline
\texttt{bus\_results} & Per-bus three-phase voltages (\texttt{ThreePhaseBusResult}) \\
\texttt{converged, iterations, residual} & Convergence indicators \\
\texttt{total\_p\_loss\_mw} & Total three-phase active power loss \\
\texttt{total\_q\_loss\_mvar} & Total three-phase reactive power loss \\
\texttt{max\_vuf\_percent} & Worst-case Voltage Unbalance Factor (\%) \\
\hline\hline
\end{tabular}
\end{center}

\texttt{ThreePhaseFlowOptions} is a lightweight wrapper over
\texttt{ThreePhaseNROptions}: \texttt{max\_iter}, \texttt{tol},
\texttt{verbose}, and \texttt{include\_shunts} are forwarded to
\texttt{analysis::solve\_three\_phase\_nr}.  No backward-forward sweep selector
is exposed.

% ------------------------------------------------------------
\subsection{PV Array Power Curve Model}
\label{sec:pv-curve}
% ------------------------------------------------------------

Declared in \texttt{power\_flow/pv\_power\_curve.hpp}.  Implements the
empirical five-parameter PV module model:

\begin{equation}
  I(V) = I_{sc}\,\left(1 - \left(\frac{V}{V_{oc}}\right)^a\right)^b
           \left(1 - c\left(\frac{V}{V_{mpp}}\right)^2\right),
\end{equation}

with fixed constants $a = 10.0$, $b = 0.5476$, $c = 0.02381$ retained for
legacy DistributionPowerFlow parity.

\texttt{compute\_pv\_power\_mw(pv)} evaluates the complete PV system output:
\begin{equation}
  P_{array} = I(V_{mpp}) \cdot V_{mpp} \cdot N_{series} \cdot N_{parallel}
              \cdot \frac{G}{1000} \cdot \eta_{inv},
\end{equation}
where $G$ is irradiance (W/m\textsuperscript{2}) and $\eta_{inv}$ is inverter
efficiency.  Falls back to \texttt{pv.p\_mw} when array parameters are absent.

% ------------------------------------------------------------
\subsection{Reliability Assessment}
\label{sec:reliability}
% ------------------------------------------------------------

Declared in \texttt{include/hacdcpf/reliability/reliability\_assessment.hpp},
implemented in \texttt{src/reliability/reliability\_assessment.cpp}.  Provides
four complementary methods for bulk-power and distribution reliability
assessment, all in namespace \texttt{hacdcpf::analysis}:

\begin{itemize}
  \item \textbf{Non-Sequential Monte Carlo} (\texttt{run\_nonsequential\_mc}) ---
        State sampling via Bernoulli trials; each sample draws an independent
        system state and evaluates load curtailment with DC-OPF.  Converges by
        Coefficient of Variation (CoV).
  \item \textbf{Sequential Monte Carlo} (\texttt{run\_sequential\_mc}) ---
        Chronological simulation driven by an hourly load profile; advances an
        event queue of exponential TTF/TTR variates to capture time correlations
        and produce annual per-year statistics.
  \item \textbf{Frequency \& Duration} (\texttt{run\_frequency\_duration\_analysis}) ---
        Analytical COPT-based method; builds a Capacity Outage Probability Table
        (COPT) by recursive convolution and computes LOLP/LOLE/LOLF without
        sampling variance.
  \item \textbf{Distribution FMEA} (\texttt{run\_distribution\_fmea}) ---
        Deterministic $N-1$ failure-mode enumeration for distribution systems;
        evaluates switching and repair stages via DC-OPF and aggregates
        frequency-weighted reliability indices.
\end{itemize}

Two utility functions populate component-level reliability data:
\texttt{apply\_ieee24\_reliability\_data(sys)} (IEEE RTS-24 generators and
branches) and \texttt{apply\_comprehensive\_reliability\_data(sys)} (full
hybrid AC/DC distribution topology including transformers, VSCs, DC branches,
static generators, renewables, PV, and storage).

\subsubsection{Options: \texttt{ReliabilityOptions}}

\begin{center}
\begin{tabular}{llll}
\hline\hline
\textbf{Field} & \textbf{Default} & \textbf{Units} & \textbf{Description} \\
\hline
\texttt{max\_iterations} & 10000 & --- & NSQ: max samples; SEQ: max years \\
\texttt{cov\_threshold} & 0.05 & --- & CoV convergence criterion \\
\texttt{curtail\_threshold\_mw} & 0.01 & MW & Minimum curtailment to count \\
\texttt{seed} & 0 & --- & RNG seed (0 = non-deterministic) \\
\texttt{load\_scale\_factor} & 1.0 & --- & Global load scaling \\
\texttt{hours\_per\_year} & 8736 & hr & Calendar (IEEE RTS uses 8736) \\
\texttt{compute\_tail\_risk} & true & --- & Compute VaR/CVaR metrics \\
\texttt{var\_confidence} & 0.95 & --- & Confidence level for VaR/CVaR \\
\texttt{compute\_distribution\_indices} & false & --- & Compute SAIFI/SAIDI/ASAI \\
\texttt{use\_importance\_sampling} & false & --- & Variance reduction (experimental) \\
\texttt{verbose} & false & --- & Print diagnostic output \\
\hline\hline
\end{tabular}
\end{center}

\subsubsection{Results: \texttt{ReliabilityResult}}

\begin{center}
\begin{tabular}{lll}
\hline\hline
\textbf{Field} & \textbf{Units} & \textbf{Description} \\
\hline
\texttt{converged} & --- & Whether CoV criterion was satisfied \\
\texttt{iterations\_used} & --- & Samples (NSQ) or years (SEQ) used \\
\texttt{final\_cov} & --- & Final Coefficient of Variation \\
\texttt{eens\_mwh\_yr} & MWh/yr & Expected Energy Not Supplied \\
\texttt{edns\_mw} & MW & Expected Demand Not Supplied \\
\texttt{lole\_hr\_yr} & hr/yr & Loss of Load Expectation \\
\texttt{lolf\_occ\_yr} & occ/yr & Loss of Load Frequency (SEQ only) \\
\texttt{plc} & --- & Probability of Load Curtailment \\
\texttt{nodal\_eens\_mwh\_yr} & MWh/yr & Per-bus EENS vector \\
\texttt{critical\_components} & --- & Sorted component importance list \\
\texttt{annual\_eens / \_lole / \_lolf} & MWh, hr, occ & Per-year values (SEQ only) \\
\texttt{tail\_risk} & --- & \texttt{TailRiskMetrics} (VaR/CVaR) \\
\texttt{distribution\_idx} & --- & \texttt{DistributionIndices} (if enabled) \\
\hline\hline
\end{tabular}
\end{center}

\subsubsection{Advanced Result Structs}

\paragraph{TailRiskMetrics.}
Computed when \texttt{compute\_tail\_risk = true}:

\begin{center}
\begin{tabular}{ll}
\hline\hline
\textbf{Field} & \textbf{Description} \\
\hline
\texttt{eens\_var}, \texttt{lole\_var} & VaR of EENS/LOLE at confidence level \\
\texttt{eens\_cvar}, \texttt{lole\_cvar} & CVaR (Expected Shortfall) of EENS/LOLE \\
\texttt{eens\_percentiles}, \texttt{lole\_percentiles} & 5th/25th/50th/75th/95th percentiles \\
\hline\hline
\end{tabular}
\end{center}

\paragraph{DistributionIndices (IEEE Std 1366-2012).}
Computed when \texttt{compute\_distribution\_indices = true} (requires
\texttt{Load::num\_customers} set):

\begin{center}
\begin{tabular}{lll}
\hline\hline
\textbf{Field} & \textbf{Units} & \textbf{Description} \\
\hline
\texttt{saifi} & int/cust/yr & System Average Interruption Frequency Index \\
\texttt{saidi} & hr/cust/yr & System Average Interruption Duration Index \\
\texttt{caidi} & hr/int & Customer Average Interruption Duration Index \\
\texttt{asai} & fraction & Average Service Availability Index \\
\texttt{asui} & fraction & Average Service Unavailability Index ($= 1 - \mathrm{ASAI}$) \\
\hline\hline
\end{tabular}
\end{center}

\paragraph{FrequencyDurationResult.}
Returned by \texttt{run\_frequency\_duration\_analysis}:

\begin{center}
\begin{tabular}{lll}
\hline\hline
\textbf{Field} & \textbf{Units} & \textbf{Description} \\
\hline
\texttt{capacity\_outage\_levels} & MW & COPT capacity outage steps \\
\texttt{cumulative\_probability} & --- & $P(\text{outage} \ge X)$ per step \\
\texttt{cumulative\_frequency} & occ/yr & $F(\text{outage} \ge X)$ per step \\
\texttt{lolp} & --- & Loss of Load Probability \\
\texttt{lole\_fd} & hr/yr & LOLE from F\&D method \\
\texttt{lolf\_fd} & occ/yr & LOLF from F\&D method \\
\texttt{lold} & hr/occ & Loss of Load Duration ($= \mathrm{LOLE}/\mathrm{LOLF}$) \\
\hline\hline
\end{tabular}
\end{center}

\subsubsection{FMEA Options and Results}

\paragraph{Key FMEAOptions fields:}

\begin{center}
\begin{tabular}{llll}
\hline\hline
\textbf{Field} & \textbf{Default} & \textbf{Units} & \textbf{Description} \\
\hline
\texttt{load\_scale\_factor} & 1.0 & --- & Global load scaling \\
\texttt{switching\_time\_hr} & 0.5 & hr & Default switching/isolation time \\
\texttt{voll} & 10000.0 & \$/MWh & Value of Lost Load \\
\texttt{enable\_microgrid\_islanding} & true & --- & Enables source/island anchor handling before physical scoring \\
\texttt{enable\_repair\_reconfiguration} & true & --- & Enables repair-stage topology candidate evaluation \\
\texttt{enable\_switch\_reconfiguration} & true & --- & Enables the implemented AC switch/tie repair search \\
\texttt{max\_repair\_switch\_actions} & 2 & --- & Max simultaneous switch actions in repair stage \\
\texttt{enable\_storage\_dispatch} & true & --- & Storage contributes bounded source capacity in the FMEA physical evaluator \\
\texttt{enable\_grid\_forming\_vsc\_support} & true & --- & Grid-forming VSCs anchor islands; VSC active transfers are dispatched in hybrid FMEA LPs \\
\hline\hline
\end{tabular}
\end{center}

\paragraph{FMEAResult system indices:}

\begin{center}
\begin{tabular}{lll}
\hline\hline
\textbf{Field} & \textbf{Units} & \textbf{Description} \\
\hline
\texttt{n\_contingencies} & --- & Total $N-1$ contingencies evaluated \\
\texttt{n\_loss\_contingencies} & --- & Contingencies causing any load loss \\
\texttt{eens\_mwh\_yr} & MWh/yr & Frequency-weighted EENS \\
\texttt{edns\_mw} & MW & Expected Demand Not Supplied \\
\texttt{lole\_hr\_yr} & hr/yr & Loss of Load Expectation \\
\texttt{lolf\_occ\_yr} & occ/yr & Loss of Load Frequency \\
\texttt{distribution\_idx} & --- & IEEE 1366 indices (SAIFI/SAIDI/ASAI) \\
\texttt{contingencies} & --- & Per-contingency detail, sorted by EENS$\downarrow$ \\
\texttt{model\_scope} & --- & \texttt{ac-only-dcopf} for AC-only systems; \texttt{hybrid-acdc-network-lp} when DC/VSC components are present \\
\texttt{model\_limitations} & --- & Mandatory report text describing linearized scope and nonlinear AC voltage/reactive limits \\
\texttt{validity} & --- & Flags for DC-load curtailment, VSC/DC power flow, AC OPF scope, nonlinear AC feasibility, and repair-stage AC/DC reconfiguration coverage \\
\hline\hline
\end{tabular}
\end{center}

Engineering reports and JSON exports that include FMEA results must surface
\texttt{model\_scope}, \texttt{model\_limitations}, and \texttt{validity}.  The
catalog enumerates AC, DC, VSC, DCDC, switch, breaker, storage, PV, and static
generator contingencies.  AC-only systems use the established DCOPF evaluator;
hybrid systems use a linear AC/DC network LP that co-optimizes AC load shedding,
DC load shedding, DC sources, DC branches, DC circuit-breaker edges, VSC
active-power transfers, and DC/DC converter transfers.  The hybrid LP does not
claim full nonlinear AC voltage/reactive feasibility.  Repair-stage explicit
switch search enumerates AC switches and normally-open AC branches only; DC
breakers, DC branches, DCDC devices, and VSC topology actions are disclosed as
not modelled by \texttt{repair\_dc\_side\_reconfiguration\_modelled=false}.

Each entry in \texttt{contingencies} (\texttt{FMEAContingencyDetail}) records
\texttt{failure\_rate}, switching-stage shed/ENS (\texttt{shed\_sw\_mw},
\texttt{ens\_sw\_mwh}), repair-stage shed/ENS (\texttt{shed\_rep\_mw},
\texttt{ens\_rep\_mwh}), annual contributions
(\texttt{eens\_contribution} $= \lambda_k(E^{sw}_k + E^{rep}_k)$,
\texttt{lole\_contribution}), per-bus nodal shed vectors, and the explicit
repair-stage switch actions selected by the topology search.

\subsubsection{Mathematical Formulation}

\paragraph{Component Failure Model.}
Each component $k$ is parameterised by annual failure rate $\lambda_k$
(occ/yr) and mean time to repair $r_k$ (hr).  For generators the forced
outage rate $\mathrm{FOR}_k$ is the primary input:
\begin{equation}
  \lambda_k = \frac{\mathrm{FOR}_k}{(1 - \mathrm{FOR}_k)\,r_k}, \qquad
  \mu_k = \frac{1}{r_k}.
\end{equation}

\paragraph{Non-Sequential Monte Carlo.}
Each sample $s$ draws an independent system state:
\begin{equation}
  X_k^{(s)} \sim \mathrm{Bernoulli}\!\left(\mathrm{FOR}_k\right), \quad
  k = 1,\ldots,K.
\end{equation}
Available capacity is evaluated by DC-OPF with load-shedding variables.
Convergence is declared when the Coefficient of Variation of the running EENS
estimate falls below \texttt{cov\_threshold}:
\begin{equation}
  \mathrm{CoV}_n = \frac{\hat{\sigma}_n}{\hat{\mu}_n \sqrt{n}} \le \epsilon_{CoV}.
\end{equation}
System indices are estimated as sample means:
\begin{align}
  \widehat{\mathrm{EENS}} &= \frac{1}{n}\sum_{s=1}^{n} C_s
    \quad [\text{MWh/yr}], \\
  \widehat{\mathrm{LOLE}} &= \frac{T_{yr}}{n}\sum_{s=1}^{n} \mathbf{1}[C_s > 0]
    \quad [\text{hr/yr}], \\
  \widehat{\mathrm{PLC}} &= \frac{1}{n}\sum_{s=1}^{n} \mathbf{1}[C_s > 0],
\end{align}
where $C_s$ (MWh) is the load curtailment in sample $s$ and $T_{yr} =
\texttt{hours\_per\_year}$.

\paragraph{Sequential Monte Carlo.}
Each component $k$ cycles through exponentially distributed times to failure
and repair:
\begin{equation}
  \mathrm{TTF}_k \sim \mathrm{Exp}(\lambda_k), \qquad
  \mathrm{TTR}_k \sim \mathrm{Exp}(\mu_k).
\end{equation}
The simulator maintains a chronological event queue, applies the hourly load
profile \texttt{LoadProfile::factors}, solves DC-OPF at each curtailment
interval, and accumulates annual statistics.  Per-year vectors
\texttt{annual\_eens/lole/lolf} enable tail-risk post-processing.

\paragraph{Frequency \& Duration (COPT).}
Generator outage states are convolved recursively.  Starting from the
all-in-service state, generator $i$ with capacity $C_i$ and FOR $q_i$ is
added via:
\begin{align}
  P(X + C_i \ge x) &\leftarrow
    (1-q_i)\,P(X \ge x) + q_i\,P(X \ge x - C_i), \\
  F(X + C_i \ge x) &\leftarrow
    (1-q_i)\,F(X \ge x) + q_i\,F(X \ge x - C_i)
    + \lambda_i\,P(X \ge x) - \lambda_i\,P(X \ge x - C_i).
\end{align}
System indices follow from the final COPT evaluated at peak load $L$:
\begin{align}
  \mathrm{LOLP} &= P(C_{out} \ge C_{avail} - L), \\
  \mathrm{LOLE} &= T_{yr} \times \mathrm{LOLP} \quad [\text{hr/yr}], \\
  \mathrm{LOLF} &= F(C_{out} \ge C_{avail} - L) \quad [\text{occ/yr}], \\
  \mathrm{LOLD} &= \mathrm{LOLE} / \mathrm{LOLF} \quad [\text{hr/occ}].
\end{align}

\paragraph{FMEA Frequency Weighting.}
For each $N-1$ contingency $k$ and stage $s \in \{\mathrm{sw}, \mathrm{rep}\}$:
\begin{equation}
  E^s_k = P^{shed}_{k,s} \cdot \tau_{k,s}, \qquad
  \mathrm{EENS} = \sum_k \lambda_k (E^{sw}_k + E^{rep}_k),
\end{equation}
\begin{equation}
  \mathrm{LOLF} = \sum_{k:\, E^{sw}_k + E^{rep}_k > 0} \lambda_k, \qquad
  \mathrm{LOLE} = \sum_k \lambda_k \!\sum_{s} \tau_{k,s}\,\mathbf{1}[P^{shed}_{k,s} > 0].
\end{equation}

\paragraph{Tail Risk (VaR/CVaR).}
Given $n$ per-sample EENS values $\{e_1,\ldots,e_n\}$ sorted ascending, the
$\alpha$-VaR and $\alpha$-CVaR at confidence level $\alpha$ are:
\begin{align}
  \mathrm{VaR}_\alpha &= e_{\lceil \alpha n \rceil}, \\
  \mathrm{CVaR}_\alpha &= \frac{1}{n(1-\alpha)}
    \sum_{i=\lceil \alpha n \rceil}^{n} e_i.
\end{align}

\subsubsection{Reference Benchmarks}

Table~\ref{tab:reliability-benchmarks} summarises expected results on the two
standard test cases.

\begin{table}[ht]
\centering
\caption{Reliability module reference benchmarks.}
\label{tab:reliability-benchmarks}
\begin{tabular}{llllll}
\hline\hline
\textbf{Test system} & \textbf{Method} & \textbf{Peak load} & \textbf{Reserve} &
  \textbf{EENS (MWh/yr)} & \textbf{LOLE (hr/yr)} \\
\hline
IEEE RTS-24 & NSQ MC & 2\,850 MW & 75\,\% & $\approx 0$ & $\approx 0$ \\
IEEE RTS-24 (stressed) & NSQ MC & 4\,600 MW & 8\,\% & $>0$ & $>0$ \\
IEEE RTS-24 & F\&D & 2\,850 MW & 75\,\% & ---\,(analytical) & $\ll 1$ \\
Comprehensive HACDCSS & FMEA & nominal & --- & $\ge 0$ & $\ge 0$ \\
Comprehensive HACDCSS (2$\times$ load) & FMEA & $2\times$ nominal & --- & $>0$ & $>0$ \\
\hline\hline
\end{tabular}
\end{table}

For the IEEE RTS-24 case with its high installed reserve ($\approx 75\,\%$
above peak), single-component contingencies rarely cause load shedding and
both EENS and LOLE converge to zero under the standard load.  The stressed
scenario (\texttt{load\_scale\_factor} $\approx 1.61$, corresponding to an
8\,\% reserve margin) produces non-trivial reliability indices consistent with
published RTS values (LOLE $\approx 9$ hr/yr, EENS $\approx 1{,}100$ MWh/yr).

% ------------------------------------------------------------
\subsection{Resilience Assessment}
\label{sec:resilience}
% ------------------------------------------------------------

Declared in \texttt{include/hacdcpf/resilience/resilience\_assess\-ment.hpp},
implemented in \texttt{src/resilience/resilience\_assessment.cpp}.

The module evaluates how well a distribution system survives and recovers
from branch outages over a configurable multi-hour horizon.  It supports
priority-weighted load shedding, network reconfiguration via tie-switch
closure, stationary storage dispatch, and Mobile Energy Storage System
(MESS) routing.

% ----------------------------------------------------------------
\subsubsection{Mathematical Model}
\label{sec:resilience-math}

\paragraph{System state at each time step.}
Let $\mathcal{B}$ be the set of buses and $\mathcal{E}$ the set of branches.
At time step $t$ with step size $\Delta t$, each branch $e\in\mathcal{E}$ has
closed/open status $z_{e,t}\in\{0,1\}$, determined by the nominal topology,
active faults, and tie-switch reconfiguration decisions.

\paragraph{Island decomposition.}
The closed subgraph $G_t = (\mathcal{B}, \{e : z_{e,t}=1\})$ is partitioned
into connected components (islands) $\mathcal{I}_1^{(t)},\ldots$,
$\mathcal{I}_m^{(t)}$ via depth-first search.  An island is \emph{grid-connected}
if it contains at least one slack or external-grid bus.

\paragraph{Island power balance.}
For each island $\mathcal{I}_k^{(t)}$, define the available supply:
\begin{equation}
  S_k^{(t)} =
    \sum_{g\in\mathcal{G}_k} P_g^{\max}
    + \sum_{r\in\mathcal{R}_k} \mu_t^{\text{res}}\,P_r^{\text{rated}}
    + \sum_{s\in\mathcal{F}_k} d_{s,t}
    + \sum_{m\in\mathcal{M}_k} d_{m,t}^{\text{MESS}},
  \label{eq:res-supply}
\end{equation}
where $\mu_t^{\text{res}}$ is the renewable availability multiplier at hour $t$,
$d_{s,t}$ is the discharge power of fixed storage unit $s$, and
$d_{m,t}^{\text{MESS}}$ is the dispatch power of MESS unit $m$ currently
serving island $k$.  If $\mathcal{I}_k^{(t)}$ is grid-connected, the grid
supplies all demand ($S_k^{(t)} = \infty$).

The total demand in island $k$ at time $t$ is:
\begin{equation}
  D_k^{(t)} = \mu_t^{\text{load}} \sum_{i\in\mathcal{L}_k} P_i^{\text{load}},
  \label{eq:res-demand}
\end{equation}
where $\mu_t^{\text{load}}$ is sampled from a 24-hour load profile.

\paragraph{Priority-weighted load shedding.}
Each load $i$ carries an importance weight $w_i \geq 1$.
The served power is allocated greedily in descending order of $w_i$:
\begin{align}
  \text{served}_i^{(t)} &= \min\!\bigl(P_i^{(t)},\; A_{\text{rem}}\bigr),
    \label{eq:res-served}\\
  A_{\text{rem}} &\leftarrow A_{\text{rem}} - \text{served}_i^{(t)},
\end{align}
with $A_{\text{rem}} = \min(D_k^{(t)}, S_k^{(t)})$ initially.
The weighted shed power and priority-tier shedding are:
\begin{equation}
  \widetilde{P}_k^{\text{shed},(t)}
    = \sum_i w_i\bigl(P_i^{(t)} - \text{served}_i^{(t)}\bigr).
  \label{eq:res-weighted-shed}
\end{equation}

\paragraph{Storage depletion.}
After dispatching, the energy state of fixed storage $s$ is updated as:
\begin{equation}
  E_{s,t+1} = \max\!\bigl(E_s^{\min},\;
    E_{s,t} - d_{s,t}\,\Delta t\,/\,\eta_s^{\text{dis}}\bigr).
  \label{eq:res-storage-soc}
\end{equation}

\paragraph{MESS travel and energy consumption.}
MESS unit $m$ at bus $b$ is routed via a Dijkstra shortest-path on the
road-network graph $G^{\text{road}}$.  Travel over route leg $\ell$ of
distance $\ell_\ell$ km consumes:
\begin{equation}
  \Delta E_m^\ell = \ell_\ell\;\kappa_m,
  \label{eq:res-mess-energy}
\end{equation}
where $\kappa_m$ is the MESS energy consumption per km (MWh/km).
Travel time for a leg is $\ell_\ell / v$ where $v$ is the travel speed.

\paragraph{Resilience index (RI).}
The primary outcome metric is the ratio of energy served to energy demanded:
\begin{equation}
  \mathrm{RI} = \frac{\sum_t \text{served}_t\,\Delta t}
                     {\sum_t D_t\,\Delta t}
             = \frac{E^{\text{served}}}{E^{\text{demand}}},
  \label{eq:resilience-index}
\end{equation}
where $\mathrm{RI}\in[0,1]$ with $\mathrm{RI}=1$ meaning no shedding occurred.

The total Energy Not Served (ENS) is:
\begin{equation}
  \mathrm{ENS} = \sum_t \bigl(D_t - \text{served}_t\bigr)\,\Delta t
              = E^{\text{demand}} - E^{\text{served}} \quad [\text{MWh}].
  \label{eq:res-ens}
\end{equation}

% ----------------------------------------------------------------
\subsubsection{Algorithm}
\label{sec:resilience-alg}

\begin{algorithm}[H]
\caption{Multi-Hour Distribution Resilience Assessment}
\label{alg:resilience}
\begin{algorithmic}[1]
\Require System $\mathcal{S}$, options \texttt{opts}
\Ensure Resilience result $\mathbf{R}$
\State Validate horizon $H > 0$ and step $\Delta t > 0$; abort if invalid
\State Build bus/branch position maps; collect base loads, grid sources,
       fixed storage, MESS, faults, transport graph
\State Expose fault sequence $\mathbf{R}.\texttt{fault\_sequence}$
\State $\mathbf{B}_0 \leftarrow$ initial branch closed status from \texttt{in\_service}
\For{$t = 0, 1, \ldots, H/\Delta t - 1$}
  \State $\mu^{\text{load}} \leftarrow \text{sampleProfile}(\text{load\_profile}, t\cdot\Delta t)$
  \State $\mu^{\text{res}} \leftarrow \text{sampleProfile}(\text{res\_profile}, t\cdot\Delta t)$
  \State Scale loads: $\mathbf{D}_t \leftarrow \mu^{\text{load}}\cdot\mathbf{D}_{\text{base}}$
  \State Advance all MESS units in transit (discrete leg stepping)
  \State $\mathbf{B}_t \leftarrow \text{applyFaults}(\mathbf{B}_{t-1}, \text{faults}, t)$
  \If{\texttt{allow\_reconfiguration}}
    \State Greedily close tie switches: prioritise connecting isolated
           demand to grid source (highest-demand island first)
  \EndIf
  \State $\{G_t^k\} \leftarrow \text{DFS-islands}(\mathbf{B}_t)$
  \For{each island $\mathcal{I}_k$}
    \State Compute $S_k$ from generators, RES, storage, MESS
    \State Allocate supply greedily by importance; compute shed and ENS
  \EndFor
  \If{\texttt{allow\_mess\_dispatch} \textbf{and} any island sheds load}
    \State Route idle MESS units to highest-deficit islands (Dijkstra + cache)
  \EndIf
  \State Update storage SOC (fixed and MESS) via~\eqref{eq:res-storage-soc}
  \If{\texttt{run\_power\_flow}}
    \State Solve adaptive AC power flow; store bus voltages and branch flows
  \EndIf
  \State Accumulate $E^{\text{demand}}, E^{\text{served}}, E^{\text{shed}}$
  \State Call progress callback; break if cancelled
  \State $\mathbf{B}_{t} \leftarrow \mathbf{B}_t$ (carry forward closed state)
\EndFor
\State $\mathrm{RI} \leftarrow E^{\text{served}} / E^{\text{demand}}$ (\eqref{eq:resilience-index})
\State Compute avg/final restoration ratio, peak shed
\State Sync legacy fields; set status to ``Completed''
\State \Return $\mathbf{R}$
\end{algorithmic}
\end{algorithm}

% ----------------------------------------------------------------
\subsubsection{Algorithm Flow Diagram}
\label{sec:resilience-flow}

\begin{center}
\begin{tikzpicture}[
  node distance = 6mm,
  proc/.style  = {draw, rounded corners=2pt, fill=blue!8,
                  text width=70mm, align=center,
                  minimum height=7mm, font=\small},
  dec/.style   = {draw, fill=orange!12, text width=64mm,
                  align=center, minimum height=7mm, font=\small,
                  rounded corners=2pt},
  arr/.style   = {-Stealth, thick},
  term/.style  = {draw, rounded corners=8pt, fill=gray!20,
                  text width=36mm, align=center,
                  minimum height=7mm, font=\small\bfseries},
]
\node[term]                    (A) {Start};
\node[proc, below=of A]        (B) {Validate options; build bus/branch maps,\\
                                    loads, faults, storage, transport graph};
\node[proc, below=of B]        (C) {Step $t$: sample profiles; advance MESS in transit};
\node[proc, below=of C]        (D) {Apply active faults and repairs to branch status};
\node[dec,  below=of D]        (E) {\textbf{Reconfig?} \quad Yes $\Rightarrow$ greedy tie-switch closure};
\node[proc, below=of E]        (F) {DFS island decomposition};
\node[proc, below=of F]        (G) {Evaluate islands: supply, priority shed, storage dispatch};
\node[dec,  below=of G]        (H) {\textbf{MESS?} \quad Yes $\Rightarrow$ route idle MESS via Dijkstra};
\node[proc, below=of H]        (I) {Update storage SOC; accumulate ENS, demand, served};
\node[dec,  below=of I]        (J) {\textbf{Power flow?} \quad Yes $\Rightarrow$ adaptive AC PF};
\node[dec,  below=of J]        (K) {\textbf{More steps \& not cancelled?}};
\node[proc, below=of K]        (L) {Compute RI, avg/final ratio, peak shed; sync fields};
\node[term, below=of L]        (M) {End};

\draw[arr] (A)--(B);
\draw[arr] (B)--(C);
\draw[arr] (C)--(D);
\draw[arr] (D)--(E);
\draw[arr] (E)--(F);
\draw[arr] (F)--(G);
\draw[arr] (G)--(H);
\draw[arr] (H)--(I);
\draw[arr] (I)--(J);
\draw[arr] (J)--(K);
\draw[arr] (K)--node[right,font=\scriptsize]{No}(L);
\draw[arr] (L)--(M);

% Loop back: Yes on K goes back to C
\draw[arr] (K.east) -- node[above,font=\scriptsize]{Yes}
  ++(8mm,0) -- ++(0,{5*(6mm+7mm)}) -| (C.north);

\end{tikzpicture}
\end{center}

% ----------------------------------------------------------------
\subsubsection{Option and Result Tables}
\label{sec:resilience-tables}

\begin{center}
\begin{tabular}{lllp{50mm}}
\hline\hline
\textbf{Field (\texttt{DistributionResilienceOptions})} & \textbf{Default} & \textbf{Units} & \textbf{Description} \\
\hline
\texttt{horizon\_hours}           & 48     & hr    & Total simulation horizon \\
\texttt{time\_step\_hr}           & 1.0    & hr    & Step size $\Delta t$ \\
\texttt{load\_scale\_factor}      & 1.0    & ---   & Global load multiplier \\
\texttt{allow\_reconfiguration}   & true   & ---   & Enable tie-switch closure \\
\texttt{allow\_mess\_dispatch}    & true   & ---   & Enable MESS dispatch \\
\texttt{max\_reconfig\_iterations}& 50     & ---   & Greedy reconfig safety limit \\
\texttt{mess\_travel\_speed\_kmph}& 40.0   & km/hr & MESS road speed \\
\texttt{default\_fault\_count}    & 2      & ---   & Auto-fault branch count \\
\texttt{default\_repair\_time\_hr}& 6.0    & hr    & Auto-fault repair duration \\
\texttt{run\_power\_flow}         & false  & ---   & Full AC PF validation \\
\hline\hline
\end{tabular}
\end{center}

\subsubsection{Strict Multi-Period MIP Path}
\label{sec:resilience-mip}

The heuristic assessment above is not the only resilience entry point.  When
\texttt{DistributionResilienceOptions::model} is
\texttt{MultiPeriodMIPLinDistFlow}, or when callers invoke
\texttt{run\_distribution\_resilience\_mip\_assessment()} directly, the code
uses \texttt{src/resilience/resilience\_restoration\_mip.cpp}.  This path builds
an independent AC-only multi-period restoration MILP instead of delegating to
the heuristic loop.

The MIP skeleton includes multi-period branch status, forest/radiality logic,
active-power LinDistFlow voltage envelopes, load-served/shedding variables,
fixed-storage energy dynamics, and time-space MESS routing arcs.  Its declared
model scope is \texttt{ac-only-lindistflow}; DC buses, DC branches, VSC
converters, and DC loads are not represented by decision variables.  The result
populates \texttt{model\_stats} with model-built/solved flags, solver name and
status, objective value, MIP gap, runtime, formulation notes,
\texttt{model\_scope}, and validity flags such as
\texttt{mip\_gap\_within\_tolerance}.

\begin{center}
\begin{tabular}{lp{75mm}}
\hline\hline
\textbf{Field (\texttt{DistributionResilienceResult})} & \textbf{Description} \\
\hline
\texttt{total\_demand\_mwh}      & Total energy demanded over horizon (MWh) \\
\texttt{total\_served\_mwh}      & Total energy served (MWh) \\
\texttt{total\_shed\_mwh}        & Total ENS (MWh); alias \texttt{total\_ens\_mwh} \\
\texttt{resilience\_index}       & RI $\in[0,1]$, see~\eqref{eq:resilience-index} \\
\texttt{avg\_restoration\_ratio} & Mean per-step restoration ratio \\
\texttt{final\_restoration\_ratio} & Restoration ratio at last step \\
\texttt{peak\_shed\_mw}          & Maximum instantaneous shed (MW) \\
\texttt{mess\_energy\_delivered\_mwh} & Total MESS dispatch energy (MWh) \\
\texttt{mess\_travel\_distance\_km}   & Cumulative MESS road distance (km) \\
\texttt{total\_switch\_actions}  & Total topology changes across all steps \\
\texttt{steps}                   & Per-step \texttt{DistributionResilienceStepResult} vector \\
\texttt{fault\_sequence}         & Faults applied (auto or user-specified) \\
\texttt{feasible, status}        & Success flag and text status \\
\hline\hline
\end{tabular}
\end{center}

% ----------------------------------------------------------------
\subsubsection{Numerical Verification}
\label{sec:resilience-numerical}

The following results are representative checks from
\texttt{tests/test\_resilience\_assessment.cpp}.  The current registered suite
contains 23 CTest entries covering heuristic resilience, the strict MIP entry
point, sequential Monte Carlo regressions, and FMEA scope/repair-search flags.

\paragraph{3-bus radial, no fault.}
A 3-bus radial feeder (Bus~1 slack, Bus~2/3 each with 1\,MW load),
4-hour horizon, no fault, flat load profile ($\mu^{\text{load}}=1$):
\[
  E^{\text{demand}} = 8\;\text{MWh},\quad
  E^{\text{shed}} = 0\;\text{MWh},\quad
  \mathrm{RI} = 1.000.
\]

\paragraph{3-bus radial, branch-1 faulted for entire 6-hour horizon.}
Branch 1 (Bus~1--Bus~2) is faulted from $t=0$; no reconfiguration or MESS:
\[
  E^{\text{demand}} > 0\;\text{MWh},\quad
  E^{\text{shed}} > 0\;\text{MWh},\quad
  \mathrm{RI} < 1.000,\quad
  P^{\text{shed}}_{\max} > 0\;\text{MW}.
\]
The algorithm correctly identifies Bus~2 and Bus~3 as an isolated island
with no local generation, shedding all load at each step.

\paragraph{4-bus ring with tie switch, branch-1 faulted.}
A ring topology (Bus~1 slack, Bus~2--4 with 1\,MW each) has a normally-open
tie switch (branch~4).  With reconfiguration enabled, the greedy algorithm
closes branch~4 at step~0, restoring full supply:
\[
  E^{\text{shed}} = 0\;\text{MWh},\quad
  \mathrm{RI} = 1.000.
\]

\paragraph{Fault repair with reconfiguration.}
Branch~1 is faulted for $[0,5)$\,hr on a 10-step horizon.  After repair
at $t=5$, greedy reconfiguration closes the branch:
\[
  \text{restoration\_ratio}(t) = 1.0 \quad \forall\; t \geq 5.
\]

\paragraph{MESS dispatch.}
A MESS (2\,MW, 10\,MWh, SOC$_0=0.9$) stationed at the slack bus is deployed
to the isolated island when branch~1 is faulted.  Over a 4-hour horizon,
total shed with MESS is less than or equal to shed without MESS:
\[
  E^{\text{shed}}_{\mathrm{MESS}} \leq E^{\text{shed}}_{\mathrm{no\text{-}MESS}}.
\]

\noindent
All 23 resilience/FMEA regression entries pass in the current
\texttt{build\_rel} suite, including the hybrid AC/DC FMEA LP and DC bus
load-scaling regressions.

% ------------------------------------------------------------
\subsection{Time-Series Power Flow}
\label{sec:time-series-pf}
% ------------------------------------------------------------

Declared in \texttt{time\_series/time\_series\_pf.hpp},
implemented in \texttt{src/time\_series/time\_series\_pf.cpp}.  Provides a
three-stage UC~$\to$~OPF~$\to$~PF pipeline for multi-period simulation.

\subsubsection{DC Network Model for Unit Commitment}

The UC stage uses a lossless DC power flow approximation over the AC network.
For $n_{bus}$ buses with susceptance matrix $\mathbf{B}$ (entry
$B_{ij} = -1/x_{ij}$), the reduced system (slack bus removed) gives:
\begin{equation}
  \mathbf{B}_{red}\,\boldsymbol{\theta}_{-s} = \mathbf{p}_{-s},
  \qquad \theta_{slack} = 0,
  \label{eq:dc-pf-uc}
\end{equation}
where $p_i = P_i^{gen} - P_i^{load}$ is the net injection at bus $i$.
Branch flow: $P_{ij} = B_{ij}(\theta_i - \theta_j)$.
Thermal limit constraint: $|P_{ij,t}| \leq F_{ij}^{max}$.

For the DC grid, an analogous conductance matrix $\mathbf{G}$ is built from
DC branch resistances ($G_{ij} = -1/r_{ij}$), with the voltage-controlled
DC bus as slack: $\mathbf{G}_{red}\,\mathbf{V}_{DC,-s} = \mathbf{P}_{DC,-s}$.

\subsubsection{Unit Commitment MILP Formulation}

\paragraph{Decision variables.}
Over $T$ time steps of width $\Delta t$ (hours):
\begin{center}
\begin{tabular}{lll}
\hline\hline
\textbf{Variable} & \textbf{Type} & \textbf{Bounds} \\
\hline
$p_{g,t}$ & continuous & $[0,\;P_g^{max}]$ \\
$u_{g,t}$ & binary & $\{0,1\}$ \\
$s_{g,t}$ & continuous & $[0,\;C_g^{su}]$ \\
$p_{s,t}$ & continuous & $[P_s^{min},\;P_s^{max}]$ (AC ESS, $+$discharge) \\
$E_{s,t}$ & continuous & $[E_s^{min},\;E_s^{max}]$ \\
$p_{r,t}$ & continuous & $[0,\;\bar{p}_{r,t}]$ (renewable, $\bar{p}_{r,t}=$ profile$\times$rated) \\
$\theta_{i,t}$ & continuous & $[-\pi,\pi]$ (with DC network) \\
$p_{vsc,c,t}$ & continuous & $[P_c^{min},\;P_c^{max}]$ (with DC network) \\
\hline\hline
\end{tabular}
\end{center}

\paragraph{Objective.}
\begin{equation}
  \min \sum_{g=1}^G\sum_{t=0}^{T-1}
    \bigl(c_g^1\,p_{g,t} + c_g^0\,u_{g,t} + s_{g,t}\bigr)\Delta t,
  \label{eq:uc-obj}
\end{equation}
where $c_g^1$ (\$/MWh), $c_g^0$ (\$/h on-cost), $s_{g,t}$ captures startup cost.

\paragraph{Constraints.}
\begin{align}
  \text{Power balance:} \quad &
    \sum_g p_{g,t} + \sum_s p_{s,t} + \sum_r p_{r,t} = P_t^{load},
    \quad \forall t, \label{eq:uc-pb}\\
  \text{Generator limits:} \quad &
    P_g^{min}\,u_{g,t} \;\leq\; p_{g,t} \;\leq\; P_g^{max}\,u_{g,t},
    \quad \forall g,t, \label{eq:uc-gen}\\
  \text{Ramp up:} \quad &
    p_{g,t} - p_{g,t-1} \;\leq\; RR_g^{up}\,\Delta t\cdot 60,
    \quad \forall g,t\geq 1, \label{eq:uc-ramp-up}\\
  \text{Ramp down:} \quad &
    p_{g,t-1} - p_{g,t} \;\leq\; RR_g^{dn}\,\Delta t\cdot 60,
    \quad \forall g,t\geq 1, \label{eq:uc-ramp-dn}\\
  \text{Startup cost:} \quad &
    s_{g,t} \;\geq\; C_g^{su}\bigl(u_{g,t} - u_{g,t-1}\bigr),
    \quad \forall g,t, \label{eq:uc-su}\\
  \text{ESS SoC:} \quad &
    E_{s,t} = E_{s,t-1}(1{-}\sigma_s)
      - \frac{p_{s,t}^{dis}\,\Delta t}{\eta_s^{dis}}
      + p_{s,t}^{chg}\,\Delta t\,\eta_s^{chg},
    \quad \forall s,t, \label{eq:uc-soc}\\
  \text{Reserve:} \quad &
    \sum_g P_g^{max}\,u_{g,t} - \sum_g p_{g,t}
    \;\geq\; R_{req}\,P_t^{load},
    \quad \forall t \;\text{(if enabled)}, \label{eq:uc-reserve}\\
  \text{DC line thermal:} \quad &
    \bigl|B_{ij}(\theta_{i,t}-\theta_{j,t})\bigr|
    \;\leq\; F_{ij}^{max}, \quad \forall (i,j),t \;\text{(if enabled)}.
    \label{eq:uc-thermal}
\end{align}
In~\eqref{eq:uc-soc}, $\sigma_s$ is the self-discharge fraction per step,
$\eta_s^{dis}$ and $\eta_s^{chg}$ are round-trip efficiency components
($\eta_s^{RT} = \eta_s^{dis}\cdot\eta_s^{chg}$), and $p_{s,t}$ is split into
discharge ($p_{s,t}^{dis} \geq 0$) and charge ($p_{s,t}^{chg} \geq 0$)
components with $p_{s,t} = p_{s,t}^{dis} - p_{s,t}^{chg}$.

\paragraph{OPF--PF Cross-Validation Metrics.}
For each step~$t$ where both solvers run:
\begin{align}
  \Delta V_{m,t} &= \max_i\bigl|V_{i,t}^{OPF} - V_{i,t}^{PF}\bigr|
    \quad [\text{pu}], \\
  \Delta\theta_t &= \max_i\bigl|\theta_{i,t}^{OPF} - \theta_{i,t}^{PF}\bigr|
    \quad [\text{rad}], \\
  \Delta P_{loss,t} &= P_{loss,t}^{PF} - P_{loss,t}^{OPF}
    \quad [\text{MW}].
\end{align}
The OPF--PF tracking bands $(\pm 20\%/5\,\text{MW}$ for slack,
$\pm 2\%/1\,\text{MW}$ for non-slack) prevent large re-dispatch between
stages.

\subsubsection{Time-Series Data Model}

The data types live in the global \texttt{hacdcpf} namespace for backward
compatibility with JSON serialisation:

\begin{center}
\begin{tabular}{lll}
\hline\hline
\textbf{Type} & \textbf{Fields} & \textbf{Description} \\
\hline
\texttt{TimeSeriesProfile} & \texttt{id, name, values} & Named time profile vector (scaling factors per step) \\
\texttt{TimeSeriesData} & \texttt{num\_steps=24} & Collection of named profiles \\
 & \texttt{step\_duration\_hr=1.0} & Step width in hours \\
 & \texttt{profiles} & Vector of \texttt{TimeSeriesProfile} \\
\hline\hline
\end{tabular}
\end{center}

\subsubsection{Unit Commitment Schedule}

\texttt{UCSchedule} carries the MILP dispatch output from Stage~I:

\begin{center}
\begin{tabular}{ll}
\hline\hline
\textbf{Field} & \textbf{Description} \\
\hline
\texttt{gen\_dispatch[g][t]} & Generator dispatch (MW) \\
\texttt{gen\_commit[g][t]} & Binary commitment status \\
\texttt{ess\_dispatch[s][t]} & Storage dispatch (MW, $+$discharge) \\
\texttt{ess\_soc[s][t]} & Storage SOC $\in [0,1]$ \\
\texttt{renewable\_dispatch[r][t]} & Renewable output (MW) \\
\texttt{vsc\_dispatch[c][t]} & VSC AC-side injection (MW, with DC network) \\
\texttt{dcdc\_dispatch[dd][t]} & DC/DC converter flow (MW, with DC network) \\
\texttt{total\_cost, feasible, solver\_name} & Schedule scalar metadata \\
\hline\hline
\end{tabular}
\end{center}

\subsubsection{OPF/PF Cross-Validation}

For each step where both OPF and PF are run, \texttt{CrossValStep} captures:

\begin{center}
\begin{tabular}{ll}
\hline\hline
\textbf{Field} & \textbf{Description} \\
\hline
\texttt{max\_vm\_diff} & $\max_i|V^{OPF}_i - V^{PF}_i|$ (pu) \\
\texttt{max\_va\_diff} & $\max_i|\theta^{OPF}_i - \theta^{PF}_i|$ relative to slack (rad) \\
\texttt{opf\_loss\_mw, pf\_loss\_mw, loss\_diff\_mw} & Loss comparison \\
\texttt{opf\_converged, pf\_converged} & Per-step convergence flags \\
\hline\hline
\end{tabular}
\end{center}

\subsubsection{Pipeline Options and Result}

\begin{center}
\begin{tabular}{lll}
\hline\hline
\textbf{Field (TimeSeriesPFOptions)} & \textbf{Default} & \textbf{Description} \\
\hline
\texttt{skip\_uc} & false & Skip UC, dispatch from profiles directly \\
\texttt{enable\_network\_constraints} & false & Nodal DC network constraints in UC MILP \\
\texttt{enable\_dc\_network\_constraints} & false & Full AC-DC coupling in UC (requires network constraints) \\
\texttt{run\_opf} & true & UC $\to$ OPF $\to$ PF pipeline; false skips OPF \\
\texttt{reserve\_requirement\_fraction} & 0.0 & Spinning reserve as fraction of demand \\
\texttt{enable\_uc\_opf\_tracking\_band} & true & Allow OPF to deviate $\pm$20\% / 5~MW from UC dispatch \\
\texttt{enforce\_non\_slack\_strict\_tracking} & true & Tighten non-slack tracking to $\pm$2\% / 1~MW \\
\texttt{verbose} & false & Print diagnostic output \\
\hline\hline
\end{tabular}
\end{center}

\begin{center}
\begin{tabular}{ll}
\hline\hline
\textbf{Field (TimeSeriesPFResult)} & \textbf{Description} \\
\hline
\texttt{opf\_results[t]} & Per-step \texttt{ACOPFResult} (if OPF enabled) \\
\texttt{num\_opf\_converged} & Count of converged OPF steps \\
\texttt{pf\_results[t]} & Per-step \texttt{PowerFlowResult} \\
\texttt{num\_converged} & Count of converged PF steps \\
\texttt{crossval[t]} & Per-step OPF/PF cross-validation \\
\texttt{num\_steps} & Total time steps simulated \\
\texttt{total\_generation\_cost} & Sum of OPF objective across all steps \\
\texttt{uc\_solve\_sec} & Wall-clock time for UC MILP (seconds) \\
\hline\hline
\end{tabular}
\end{center}

\subsubsection{Public API}

\begin{center}
\begin{tabular}{ll}
\hline\hline
\textbf{Function} & \textbf{Description} \\
\hline
\texttt{solve\_unit\_commitment(sys, ts, opts)} & Run UC MILP only; returns \texttt{UCSchedule} \\
\texttt{solve\_time\_series\_pf(sys, ts, opts)} & Full UC $\to$ OPF $\to$ PF pipeline; returns \texttt{TimeSeriesPFResult} \\
\hline\hline
\end{tabular}
\end{center}

% ------------------------------------------------------------
\subsection{Annual Production Simulation}
\label{sec:annual-production-sim}
% ------------------------------------------------------------

Declared in \texttt{time\_series/annual\_production\_sim.hpp},
implemented in \texttt{src/time\_series/annual\_production\_sim.cpp}.
Ported from the reference implementation in Sipeng Luo's
\textit{HybridACDCPowerSystemsPlanning} repository.

\texttt{solve\_annual\_production\_simulation(sys, ts, opts)} solves a full
year of system operation through a four-level hierarchical temporal
decomposition:

\begin{enumerate}
  \item \textbf{L0 --- Annual planning}: Coarse monthly (or weekly)
        maintenance scheduling with aggregate energy and fuel budgets, storage
        SoC boundary conditions, and block-level cost targets.
  \item \textbf{L1 --- Monthly coordination}: Bridges annual plan to weekly
        UC solves; enforces cyclic SoC continuity across block boundaries.
  \item \textbf{L2 --- Weekly unit commitment}: 168-step binding MILP window
        with 48-hour lookahead; inherits energy budgets from L0/L1.
  \item \textbf{L3 --- Daily UC $\to$ OPF $\to$ PF replay}: 24-step daily
        sub-horizon; pipes each UC schedule into the full
        \texttt{solve\_time\_series\_pf} OPF/PF pipeline.
\end{enumerate}

\subsubsection{Mathematical Formulation}

\paragraph{Block Range Construction.}
Given $T_{yr}$ total time steps of width $\Delta t$ (hours):
\begin{equation}
  T_{yr} = \lfloor 8760 / \Delta t \rfloor.
\end{equation}
For \texttt{Monthly} blocks, block $b$ spans $d_b\lfloor 24/\Delta t\rfloor$
steps, where $d_b \in \{28,29,30,31\}$ (days per month; January absorbs
any rounding remainder).  For \texttt{Weekly} blocks:
$T_b = \lfloor 168/\Delta t\rfloor$ for 52 blocks.

\paragraph{L0 Annual Planning Optimisation.}
The coarse-level problem determines maintenance schedules and energy budgets
per block $b \in \{1,\ldots,N_B\}$:
\begin{equation}
  \min_{\mathbf{u}^{maint},\mathbf{E}^{budget}} \;
    \sum_{b=1}^{N_B} c_b^{plan}
    \label{eq:l0-obj}
\end{equation}
\begin{align}
  \text{s.t.} \quad &
    x_{g,b}^{maint} \in \{0,1\},
    \quad \sum_b x_{g,b}^{maint} = K_g^{maint}, \quad \forall g,
    \label{eq:l0-maint}\\
    &E_{g,b} \leq E_g^{budget,yr}, \quad \forall g,b,
    \label{eq:l0-energy}\\
    &\sum_g h_g\,E_{g,b} \leq F_b^{budget}, \quad \forall b,
    \label{eq:l0-fuel}\\
    &E_{s,0} = E_{s,T-1} \quad \text{(cyclic SoC, if \texttt{enforce\_cyclic\_soc})},
    \label{eq:cyclic-soc}
\end{align}
where $K_g^{maint}$ is the required maintenance block count for generator $g$
and $h_g$ is the heat rate.

\paragraph{L2 Weekly Rolling Horizon.}
Each week solves a rolling MILP with a binding window $[0, T_{bind})$ and
an advisory look-ahead $[T_{bind}, T_{bind}+T_{look})$:
\begin{equation}
  \min_{\mathbf{x}} \;
  \sum_{t=0}^{T_{bind}+T_{look}-1} c_t(\mathbf{x}_t)
  \label{eq:l2-obj}
\end{equation}
subject to the UC constraints~\eqref{eq:uc-gen}--\eqref{eq:uc-soc} plus the
inherited energy budget from L0:
\begin{equation}
  \sum_{t=0}^{T_{bind}-1} p_{g,t}\,\Delta t \;\leq\; E_{g,b}^{remain}.
  \label{eq:l2-budget}
\end{equation}
Only $\mathbf{x}_{0:T_{bind}-1}$ are committed; the look-ahead window is
re-solved in the next rolling step.

\paragraph{Annual Cost and Energy Aggregates.}
\begin{align}
  J^{yr}     &= \sum_{b=1}^{N_B} c_b^{total}, &
  E^{gen,yr} &= \sum_{t=0}^{T_{yr}-1} P_t^{gen}\,\Delta t,\\
  E^{RE,yr}  &= \sum_{t=0}^{T_{yr}-1} P_t^{ren}\,\Delta t, &
  E^{ENS,yr} &= \sum_{t=0}^{T_{yr}-1} P_t^{shed}\,\Delta t.
\end{align}

\subsubsection{Decomposition Data Types}

\begin{center}
\begin{tabular}{ll}
\hline\hline
\textbf{Type} & \textbf{Description} \\
\hline
\texttt{AnnualPlanBlock} & L0 block (month/week): maintenance flags, energy budgets, fuel cap, SoC boundary \\
\texttt{AnnualPlanResult} & Collection of blocks with total plan cost and feasibility \\
\texttt{WeeklySchedule} & L2 UC schedule with binding window, lookahead, and remaining energy budget \\
\hline\hline
\end{tabular}
\end{center}

\subsubsection{Per-Step and Per-Block Results}

\begin{center}
\begin{tabular}{ll}
\hline\hline
\textbf{Type / Field} & \textbf{Description} \\
\hline
\texttt{AnnualStepResult} & Lightweight per-hour snapshot: generation, load, renewable, curtailment, ESS, loss, ENS, convergence \\
\texttt{BlockSummary} & Monthly aggregation: energy MWh totals, ENS, cost, PF/OPF convergence count \\
\texttt{PFSnapshot} & Sampled voltage/branch-flow snapshot for dashboard visualisation \\
\hline\hline
\end{tabular}
\end{center}

\subsubsection{Component Annual Statistics}

\begin{center}
\begin{tabular}{ll}
\hline\hline
\textbf{Type} & \textbf{Key fields} \\
\hline
\texttt{GenAnnualStats} & Total energy MWh, capacity factor, startups/shutdowns, hours online \\
\texttt{StorageAnnualStats} & Total charge/discharge MWh, cycle count \\
\texttt{RenewableAnnualStats} & Total energy MWh, curtailed MWh, capacity factor, curtailment rate \\
\hline\hline
\end{tabular}
\end{center}

\subsubsection{Annual Production Simulation Options}

\begin{center}
\begin{tabular}{lll}
\hline\hline
\textbf{Field} & \textbf{Default} & \textbf{Description} \\
\hline
\texttt{block\_type} & Monthly & L0 decomposition granularity (Monthly or Weekly) \\
\texttt{weekly\_lookahead\_hours} & 48 & L2 advisory look-ahead buffer \\
\texttt{daily\_window\_hours} & 24 & L3 daily sub-horizon size \\
\texttt{enforce\_cyclic\_soc} & true & Require $E_{s,T-1} = E_{s,0}$ \\
\texttt{iterative\_feedback} & false & Bottom-up $\to$ top-level re-solve loop \\
\texttt{max\_feedback\_iterations} & 3 & Max feedback rounds \\
\texttt{skip\_replay} & false & Schedule-only mode (no OPF/PF) \\
\texttt{pf\_snapshot\_interval} & 24 & Store PF snapshot every $N$ steps (0 = off) \\
\texttt{curtailment\_penalty} & 50 & Curtailment penalty (\$/MWh) \\
\texttt{ens\_penalty} & 10000 & Energy not served penalty (\$/MWh) \\
\hline\hline
\end{tabular}
\end{center}

\subsubsection{Annual Production Simulation Result}

\texttt{AnnualProductionSimResult} aggregates the full hierarchy output:

\begin{center}
\begin{tabular}{ll}
\hline\hline
\textbf{Field} & \textbf{Description} \\
\hline
\texttt{annual\_plan} & L0 \texttt{AnnualPlanResult} \\
\texttt{weekly\_schedules} & Vector of L2 \texttt{WeeklySchedule} \\
\texttt{step\_results[t]} & Per-hour \texttt{AnnualStepResult} (8760 entries) \\
\texttt{pf\_snapshots} & Sampled \texttt{PFSnapshot} for visualisation \\
\texttt{monthly\_summaries} & Per-block \texttt{BlockSummary} \\
\texttt{gen\_stats, storage\_stats, renewable\_stats} & Component-level annual statistics \\
\texttt{total\_cost, total\_gen\_mwh, total\_load\_mwh} & Scalar annual aggregates \\
\texttt{total\_renewable\_mwh, total\_curtailment\_mwh} & Renewable integration metrics \\
\texttt{total\_ens\_mwh, total\_loss\_mwh} & Reliability and loss aggregates \\
\texttt{num\_pf\_converged, num\_opf\_converged} & Convergence counts \\
\texttt{feasible, solver\_name} & Overall feasibility and solver used \\
\hline\hline
\end{tabular}
\end{center}

% ------------------------------------------------------------
\subsection{Lifecycle Simulation}
\label{sec:lifecycle-sim}
% ------------------------------------------------------------

Declared in \texttt{time\_series/lifecycle\_simulation.hpp},
implemented in \texttt{src/time\_series/lifecycle\_simulation.cpp}.
Ported from the reference implementation in Sipeng Luo's
\textit{HybridACDCPowerSystemsPlanning} repository.

\texttt{run\_lifecycle\_simulation(sys, ts, opts)} runs a multi-year (default
20~year) planning horizon, advancing the system state year by year:

\begin{enumerate}
  \item For each year, scale loads by the cumulative growth factor and
        apply current component capacities (accounting for degradation and
        derating).
  \item Run the annual production simulation (Tier~1).
  \item Compute annual carbon emissions with stratified-sampling carbon
        attribution (Tier~2) and theoretical error bounds (Tier~3).
  \item Apply battery state-of-health (SoH) degradation (cycle and
        calendar) and trigger replacement when SoH falls below
        \texttt{replacement\_soh\_threshold}.
  \item Apply annual renewable derating.
  \item Discount annual costs at \texttt{discount\_rate} and accumulate
        the net present value of the total lifecycle cost.
\end{enumerate}

\subsubsection{Mathematical Formulation}

\paragraph{Demand and Supply Scaling.}
In year $y$ ($y=1,\ldots,N_Y$), the time-varying load profile is scaled by the
cumulative growth factor:
\begin{equation}
  P_t^{load}(y) = P_t^{load}(1)\cdot(1+\rho_{lg})^{y-1},
\end{equation}
and each renewable's rated capacity declines by annual derating rate
$\alpha_r^{der}$:
\begin{equation}
  P_{r,y}^{rated} = P_{r,0}^{rated}\cdot(1-\alpha_r^{der})^{y-1}.
\end{equation}

\paragraph{Battery State-of-Health (SoH) Degradation.}
Two degradation mechanisms are tracked independently.
\textit{Cycle ageing} accumulates full-equivalent cycles $N_{s,y}^{cum}$ from
the annual production simulation; \textit{calendar ageing} grows linearly with
time:
\begin{align}
  \text{SoH}_{s,y}^{cycle} &= 1 - \alpha_s^{cyc}\,N_{s,y}^{cum},
    \label{eq:soh-cyc} \\
  \text{SoH}_{s,y}^{cal}   &= 1 - \alpha_s^{cal}\,(y-1),
    \label{eq:soh-cal} \\
  \text{SoH}_{s,y}         &= \min\!\left(\text{SoH}_{s,y}^{cycle},\;
                               \text{SoH}_{s,y}^{cal}\right).
    \label{eq:soh-combined}
\end{align}
When $\text{SoH}_{s,y} \leq \text{SoH}_{EoL}$
(\texttt{replacement\_soh\_threshold}), the battery is replaced (SoH reset
to~1) and the replacement capital cost $C_s^{rep}$ is added to year~$y$.
The effective energy capacity in year~$y$ is:
\begin{equation}
  E_s^{eff}(y) = E_s^{rated}\cdot\text{SoH}_{s,y}.
\end{equation}

\paragraph{Net Present Value.}
\begin{equation}
  J^{NPV} = \sum_{y=1}^{N_Y} \frac{J_y^{ann} + C_y^{rep}}{(1+r_d)^{y-1}},
  \label{eq:npv}
\end{equation}
where $r_d$ is the annual discount rate (\texttt{discount\_rate}),
$J_y^{ann}$ is the total operational cost from the annual production simulation,
and $C_y^{rep}$ is the battery replacement cost in year~$y$ (zero if no
replacement occurs).

\subsubsection{Stratified Carbon Attribution (Tier~2)}

The full-year carbon estimate is computed from a stratified random sample of
operating hours to reduce the computational cost of per-hour power-flow
evaluations.

\paragraph{Tier-1 Dense Baseline.}
Using the mean fleet emission intensity $\bar{e}$ (tCO\textsubscript{2}/MWh):
\begin{equation}
  \hat{C}_y^{dense} = \bar{e}\sum_{t=0}^{T_{yr}-1} P_t^{gen}\,\Delta t.
  \label{eq:tier1}
\end{equation}

\paragraph{Strata Classification.}
The $T_{yr}$ hours are partitioned into $M=7$ strata according to system
operating state (implemented in \texttt{classify\_strata()}):

\begin{center}
\begin{tabular}{lll}
\hline\hline
\textbf{Stratum} & \textbf{Index} & \textbf{Classification condition} \\
\hline
Peak Load       & 0 & $P_t^{load} \geq 0.85\,P_{load}^{max}$ \\
Low Renewable   & 1 & $P_t^{ren} \leq 0.10\,P_{ren}^{max}$ \\
High Curtailment & 2 & $c_t \geq 0.30\,c_{max}$ \\
Heavy Charging  & 3 & $P_{ess,t} < -0.50\,P_{chg}^{max}$ \\
Heavy Discharging & 4 & $P_{ess,t} > +0.50\,P_{dis}^{max}$ \\
Low Load        & 5 & $P_t^{load} \leq 0.30\,P_{load}^{max}$ \\
Normal          & 6 & all remaining hours \\
\hline\hline
\end{tabular}
\end{center}

Each hour is placed in the first matching stratum (priority order as listed).
$N_m = |\mathcal{H}_m|$ is the stratum population size.

\paragraph{Neyman Sample Allocation.}
The total sample budget $n = \sum_m n_m$ is distributed proportional to
$N_m\,\hat\sigma_m$:
\begin{equation}
  n_m = \left\lfloor n\cdot
    \frac{N_m\,\hat\sigma_m}{\sum_{m'} N_{m'}\,\hat\sigma_{m'}}
  \right\rfloor,
\end{equation}
with at least one sample per non-empty stratum.

\paragraph{Stratified Estimator.}
Let $\xi_t = \bar{e}\,P_t^{gen}\,\Delta t$ and
$\bar\xi_m = n_m^{-1}\sum_{t\in\mathcal{S}_m}\xi_t$.  Then:
\begin{equation}
  \hat{C}_y^{samp} = \sum_{m=1}^M N_m\,\bar\xi_m,
  \label{eq:tier2}
\end{equation}
with stratified variance:
\begin{equation}
  \hat{V}\!\left[\hat{C}_y^{samp}\right]
  = \sum_{m=1}^M \frac{N_m^2}{n_m}\,\hat\sigma_m^2
    \!\left(1 - \frac{n_m}{N_m}\right),
  \label{eq:strat-var}
\end{equation}
where $\hat\sigma_m^2$ is the sample variance within stratum $m$.

\subsubsection{Theoretical Error Bounds (Tier~3)}

The lifecycle module computes guaranteed upper bounds on the carbon
estimation error across three error sources:

\begin{center}
\begin{tabular}{lll}
\hline\hline
\textbf{Bound type} & \textbf{Struct} & \textbf{Formula / parameter} \\
\hline
Dispatch approximation & \texttt{DispatchErrorBound} & $\varepsilon_{\text{disp}} \leq \bar{e}\sum_t L_{\max} P_{\text{load}} \Delta t$ \\
Stratified sampling & \texttt{SamplingErrorBound} & $\varepsilon_{\text{samp}} \leq z_{1-\alpha/2}\sqrt{\hat{\sigma}^2_{\text{strat}}}$ \\
Storage carbon propagation & \texttt{StorageCarbonErrorBound} & $\varepsilon_{\text{stor}} \leq K_w \Delta_{\max}$ \\
\hline\hline
\end{tabular}
\end{center}

\paragraph{Dispatch bound} (UC linearisation error):
\begin{equation}
  \varepsilon_{disp} = \bar{e}\cdot L_{max}\sum_t P_t^{load}\,\Delta t,
  \label{eq:bound-disp}
\end{equation}
where $L_{max}$ is the per-step Lipschitz constant for the dispatch-to-carbon
map.

\paragraph{Sampling bound} ($95\,\%$ confidence half-width):
\begin{equation}
  \varepsilon_{samp} = z_{1-\alpha/2}\sqrt{\hat{V}\!\left[\hat{C}_y^{samp}\right]},
  \label{eq:bound-samp}
\end{equation}
with $z_{0.975} = 1.96$ for a two-sided $\alpha = 0.05$ interval.

\paragraph{Storage carbon propagation bound}:
\begin{equation}
  \varepsilon_{stor} = K_w\cdot\Delta_{max}\cdot E_{ess}^{total},
  \label{eq:bound-stor}
\end{equation}
where $K_w = \max_t|\delta P_t^{gen}|\times 10^{-3}$ is the Lipschitz
constant from the maximum generation shift, and $\Delta_{max}$ is the
maximum temporal gap between sampled hours.

\paragraph{Combined bound}:
\begin{equation}
  \varepsilon_{total} = \varepsilon_{disp} + \varepsilon_{samp} + \varepsilon_{stor}.
  \label{eq:bound-total}
\end{equation}

Per-year combined bounds are stored in \texttt{YearlyErrorBounds}; cross-
validation between the dense Tier~1 and sampled Tier~2 estimates is
recorded in \texttt{BoundCrossValidation}.  The cross-validation metrics are:
\begin{align}
  \text{SamplingGap}    &= \frac{|\hat{C}^{dense} - \hat{C}^{samp}|}
                            {\hat{C}^{dense}}, \\
  \text{BoundTightness} &= \frac{\varepsilon_{total}}{\hat{C}^{dense}}.
\end{align}

\subsubsection{Lifecycle Simulation Options}

\begin{center}
\begin{tabular}{lll}
\hline\hline
\textbf{Field} & \textbf{Default} & \textbf{Description} \\
\hline
\texttt{num\_years} & 20 & Planning horizon \\
\texttt{discount\_rate} & 0.05 & Annual NPV discount rate \\
\texttt{load\_growth\_rate} & 0.02 & Annual load growth rate \\
\texttt{cycle\_degradation\_per\_cycle} & $4\times10^{-5}$ & SoH loss per cycle (25\,000 cycles $\to$ 100\%) \\
\texttt{calendar\_degradation\_per\_year} & 0.005 & Annual calendar SoH loss \\
\texttt{replacement\_soh\_threshold} & 0.7 & Replace battery when SoH $<$ 70\% \\
\texttt{replacement\_cost\_per\_kwh} & 250 & Replacement cost (\$/kWh) \\
\texttt{renewable\_derating\_per\_year} & 0.005 & Annual renewable capacity derating \\
\texttt{step\_duration\_hr} & 6.0 & Intra-year time step size (hours) \\
\hline\hline
\end{tabular}
\end{center}

\subsubsection{Lifecycle Result Types}

\begin{center}
\begin{tabular}{ll}
\hline\hline
\textbf{Type / Field} & \textbf{Description} \\
\hline
\texttt{YearResult} & Full annual result: cost, energy, carbon, SoH states, error bounds, replacements \\
\texttt{StorageYearState} & Per-storage: SoH (cycle + calendar), effective capacity, cycle count, replacement flag \\
\texttt{RenewableYearState} & Per-renewable: derated capacity, original capacity, derating factor \\
\texttt{ReplacementEvent} & Year, storage index/name, pre-replacement SoH, cost \\
\hline\hline
\end{tabular}
\end{center}

\begin{center}
\begin{tabular}{ll}
\hline\hline
\textbf{Field (LifecycleSimResult)} & \textbf{Description} \\
\hline
\texttt{year\_results} & Ordered \texttt{YearResult} for each year \\
\texttt{npv\_total\_cost} & Net present value of total costs over the horizon \\
\texttt{total\_carbon\_tco2} & Cumulative carbon emissions (tCO\textsubscript{2}) \\
\texttt{total\_replacement\_cost} & Sum of all battery replacement costs \\
\texttt{total\_replacements} & Count of replacement events \\
\texttt{all\_replacements} & Flat list of all \texttt{ReplacementEvent}s \\
\texttt{feasible, summary\_text} & Overall feasibility and text summary \\
\hline\hline
\end{tabular}
\end{center}

\subsubsection{Capacity Scaling and Comparison}

\texttt{apply\_capacity\_scaling(sys, scaling)} scales component capacities
in-place before a lifecycle run, enabling long-term planning studies:

\begin{center}
\begin{tabular}{ll}
\hline\hline
\textbf{Field (CapacityScaling)} & \textbf{Applies to} \\
\hline
\texttt{pv\_scale} & PV \texttt{pmax\_mw} \\
\texttt{wind\_scale} & Wind generator \texttt{p\_rated\_mw} \\
\texttt{bess\_power\_scale} & Battery \texttt{p\_rated\_mw}, \texttt{pmax}, \texttt{pmin} \\
\texttt{bess\_energy\_scale} & Battery \texttt{e\_rated\_mwh} \\
\texttt{diesel\_scale} & Diesel/gas generator \texttt{pmax} \\
\hline\hline
\end{tabular}
\end{center}

\texttt{run\_lifecycle\_comparison(sys, ts, opts, sweep\_param, scale\_values)}
sweeps one parameter (e.g.\ \texttt{"pv"}, \texttt{"bess\_power"}) across
the provided scale factors and returns a \texttt{LifecycleCompareResult}
containing a \texttt{ScenarioResult} for each point, with per-year carbon
and cost trajectories and installed capacity metadata.

% ------------------------------------------------------------
\subsection{Short-Circuit Analysis (IEC~60909)}
\label{sec:short-circuit-math}
% ------------------------------------------------------------

Declared in \texttt{include/hacdcpf/analysis/short\_circuit.hpp},
implemented in \texttt{src/analysis/short\_circuit.cpp}.
The solver computes subtransient short-circuit currents, short-circuit
power, peak current, and remaining voltages for arbitrary fault types
according to IEC~60909:2016.

\subsubsection{Voltage Factor}

IEC~60909 introduces a voltage factor $c$ to account for pre-fault
operating conditions:
\begin{center}
\begin{tabular}{lll}
\hline\hline
\textbf{Nominal voltage} & \textbf{Maximum ($c_{\max}$)} & \textbf{Minimum ($c_{\min}$)} \\
\hline
$V_n > 1\,\text{kV}$ (HV)           & 1.10 & 1.00 \\
$V_n \in (0.1,\,1]\,\text{kV}$ (LV) & 1.10 & 0.90 \\
$V_n \le 0.1\,\text{kV}$            & 1.05 & 0.95 \\
\hline\hline
\end{tabular}
\end{center}
The initial symmetrical short-circuit voltage source at the fault bus
is modelled as $\underline{V}_0 = c\,V_n$ in the equivalent circuit.

\subsubsection{Sequence Admittance Matrices}

Three symmetric admittance matrices $\mathbf{Y}_1$, $\mathbf{Y}_2$,
$\mathbf{Y}_0 \in \mathbb{C}^{n\times n}$ are assembled from:
\begin{itemize}
  \item \textbf{Branches} (lines and 2-winding transformers): standard
        $\pi$-model with series admittance $y_s = (r+jx)^{-1}$ and
        shunt susceptance $b/2$.  For off-nominal tap ratio $\tau$:
        \begin{equation}
          Y_{ii}^{(1)} \mathrel{+}= \frac{y_s}{|\tau|^2} + \frac{jb}{2|\tau|^2},
          \quad Y_{jj}^{(1)} \mathrel{+}= y_s + \frac{jb}{2},
          \quad Y_{ij}^{(1)} \mathrel{+}= -\frac{y_s}{\tau^*}.
          \label{eq:sc-ybus-branch}
        \end{equation}
        Zero-sequence branches use $(r_0, x_0, b_0)$ identically.
  \item \textbf{Synchronous generators}: subtransient impedance
        $\underline{z}_g = (r_a + jx''_d)\,(S_{\text{base}}/S_g)$ is
        corrected by the KG factor:
        \begin{equation}
          K_G = \frac{c}{1 + x''_d\,\sin\phi_G},
          \qquad
          \underline{z}_g^{\text{corr}} = K_G\,\underline{z}_g,
          \label{eq:sc-KG}
        \end{equation}
        where $\sin\phi_G = \sqrt{1-\cos^2\!\phi_G}$.
        The corrected shunt admittance $1/\underline{z}_g^{\text{corr}}$
        is added to the diagonal of $\mathbf{Y}_1$ and $\mathbf{Y}_2$
        (since $Z_2 \approx Z_1$ for synchronous machines).
  \item \textbf{Asynchronous motors}: modelled as an impedance shunt
        $\underline{z}_m = (r_m + jx_m)\,S_{\text{base}}/(V_n^2)$ added
        to $\mathbf{Y}_1$ and $\mathbf{Y}_2$.
  \item \textbf{External grids}: when the maximum short-circuit current
        $I_{kq}$ is given, the equivalent impedance is
        $z_{\text{ext}} = c\,U_{nq}/(\sqrt{3}\,I_{kq})$; otherwise
        direct $(r, x)$ shunt.
  \item \textbf{Grid-forming VSC converters}: source behind filter
        impedance $\underline{z}_{\text{conv}} = (r_1 + jx_1)
        (S_{\text{base}}/S_{\text{rated}})$ added to $\mathbf{Y}_1$.
  \item \textbf{3-winding transformers}: star-branch equivalent with
        arm impedances
        $Z_H = \tfrac{1}{2}(Z_{HM}+Z_{HL}-Z_{ML})$, etc.; Kron-reduced
        as in the 2W case.
\end{itemize}

\subsubsection{Thévenin Impedance and Initial Symmetric Current}

The Thévenin impedance at fault bus $k$ is the diagonal element of the
positive-sequence Z-bus:
\begin{equation}
  Z_{kk}^{(1)} = \bigl[\mathbf{Y}_1^{-1}\bigr]_{kk}.
  \label{eq:sc-Zkk}
\end{equation}
In the implementation $Z_{kk}$ is obtained efficiently via a single
right-hand-side solve ($\mathbf{Y}_1\,\mathbf{e}_k = \mathbf{y}$, then
$Z_{kk} = \mathbf{e}_k^\top\mathbf{y}$) using a pre-factored SparseLU
decomposition, avoiding a full matrix inversion.

\paragraph{Fault-type effective impedances.}
The effective impedance $Z_k^{\text{eff}}$ and initial symmetrical
short-circuit current $I''_k$ depend on the fault type:

\begin{center}
\renewcommand{\arraystretch}{1.4}
\begin{tabular}{lll}
\hline\hline
\textbf{Fault type} & \textbf{Effective impedance $Z_k^{\text{eff}}$}
  & \textbf{$I''_k$ (pu, base $V = 1$)} \\
\hline
3-phase (3PH)
  & $Z_1 + Z_f$
  & $c\,/\,|Z_k^{\text{eff}}|$ \\
Single-line-to-ground (SLG)
  & $(Z_1 + Z_2 + Z_0 + 3Z_f)/3$
  & $c\,/\,|Z_k^{\text{eff}}|$ \\
Line-to-line (LL)
  & $(Z_1 + Z_2)/\sqrt{3}$
  & $c\,/\,|Z_k^{\text{eff}}|$ \\
Double-line-to-ground (DLG)
  & $Z_1 + Z_2 \,\|\, Z_0 + Z_f$
  & $c\,/\,|Z_k^{\text{eff}}|$ \\
\hline\hline
\end{tabular}
\renewcommand{\arraystretch}{1.0}
\end{center}
where $Z_2 \approx Z_1$ for passive networks, $Z_f$ is the fault
impedance (zero for bolted faults), and $\,\|\,$ denotes parallel
combination.

\paragraph{Short-circuit MVA.}
\begin{equation}
  S''_k = S_{\text{base}}\,\frac{c^2}{\left|Z_{kk}^{(1)}\right|}.
  \label{eq:sc-Sk}
\end{equation}

\subsubsection{Peak Current and Kappa Factor}

The peak (asymmetric) short-circuit current accounts for the decaying
DC offset.  The peak factor $\kappa$ is defined by the equivalent
$R/X$ ratio at the fault bus:
\begin{equation}
  \kappa = 1.02 + 0.98\,e^{-3R/X}, \qquad \kappa \in [1.0,\,2.0].
  \label{eq:sc-kappa}
\end{equation}
For meshed networks with Method~B (IEC~60909 §C.3), $\kappa$ is
multiplied by 1.15 before clamping.  The peak current is:
\begin{equation}
  i_p = \kappa\,\sqrt{2}\,I''_k.
  \label{eq:sc-ip}
\end{equation}

\subsubsection{Breaking Current}

The symmetrical breaking current $I_b$ accounts for AC-component decay
during the fault-clearing interval $t_{\text{break}}$:
\begin{equation}
  I_b = \mu\,I''_k,
  \label{eq:sc-Ib}
\end{equation}
where the decay factor $\mu$ (IEC~60909 \S4.7) is:
\begin{equation}
  \mu(t_{\text{break}}) = \begin{cases}
    0.84 + 0.26\,e^{-0.26\,I''_k/I_r} & t_{\text{break}} \le 0.02\,\text{s}, \\
    0.71 + 0.51\,e^{-0.30\,I''_k/I_r} & t_{\text{break}} \le 0.05\,\text{s}, \\
    0.62 + 0.72\,e^{-0.32\,I''_k/I_r} & t_{\text{break}} \le 0.10\,\text{s}, \\
    0.56 + 0.94\,e^{-0.38\,I''_k/I_r} & \text{otherwise},
  \end{cases}
\end{equation}
with $I_r$ being the rated current.  For asynchronous motor contributions
a $q$-factor correction applies:
\begin{equation}
  I_{b,m} = \mu\,q\,I''_{km}, \qquad
  q \approx 0.26 + 0.10\ln(P_n/p), \quad p = \text{number of pole pairs}.
\end{equation}

\subsubsection{Thermal Equivalent Current}

The equivalent thermal current (for cable rating and equipment
verification) is computed per IEC~60909 \S4.9:
\begin{equation}
  I_{\text{th}} = I''_k\,\sqrt{m + n},
  \label{eq:sc-Ith}
\end{equation}
where the $m$ factor captures the DC-component heating integral and
the $n$ factor captures the AC-component decay integral over the
short-circuit duration $T_k$.

\subsubsection{Voltage at Non-Faulted Buses}

The remaining (residual) voltage at any bus~$j$ during a fault at
bus~$k$ is (IEC~60909 \S4.3):
\begin{equation}
  \hat{V}_j = c\,\left(1 - \frac{Z_{jk}^{(1)}}{Z_{kk}^{(1)}}\right),
  \label{eq:sc-Vrem}
\end{equation}
expressed in per-unit on the nominal voltage.  For unbalanced faults
the transfer impedances from all three sequence networks are combined
according to the fault-type effective-impedance formula.

\subsubsection{Options and Result Types}

\begin{center}
\begin{tabular}{lll}
\hline\hline
\textbf{Field (SCDetailedOptions)} & \textbf{Default} & \textbf{Description} \\
\hline
\texttt{c\_factor} & 1.10 & IEC~60909 voltage factor $c$ \\
\texttt{fault\_type} & \texttt{ThreePhase} & Fault type (3PH, SLG, LL, DLG) \\
\texttt{calc\_type} & \texttt{Max} & Maximum or minimum calculation \\
\texttt{kappa\_method} & \texttt{A} & Kappa method (A = radial, B = meshed $\times 1.15$) \\
\texttt{topology} & \texttt{Meshed} & Network topology hint for $\kappa$ \\
\texttt{default\_xdpp} & 0.1 & Default $x''_d$ when not specified (pu) \\
\texttt{t\_break} & 0.1\,s & Breaking time for $\mu$/$I_b$ computation \\
\hline\hline
\end{tabular}
\end{center}

\begin{center}
\begin{tabular}{ll}
\hline\hline
\textbf{Field (SCBusResult)} & \textbf{Description} \\
\hline
\texttt{bus\_id} & Bus identifier \\
\texttt{ikss\_ka} & Initial symmetrical s.c.\ current $I''_k$ (kA, total) \\
\texttt{ikss\_1\_ka} & $I''_k$ without converter contributions (kA) \\
\texttt{ikss\_no\_motor\_ka} & $I''_k$ less motor contributions (kA) \\
\texttt{ip\_ka} & Peak current $i_p$ (kA) \\
\texttt{ib\_ka} & Symmetrical breaking current $I_b$ (kA) \\
\texttt{sk\_mva} & Short-circuit power $S''_k$ (MVA) \\
\texttt{z\_thevenin} & Thévenin impedance $Z_{kk}$ (pu, complex) \\
\texttt{v\_remaining} & Residual voltage vector at all buses (pu) \\
\hline\hline
\end{tabular}
\end{center}

% ============================================================
\subsection{Three-Phase Newton-Raphson Power Flow}
\label{sec:three-phase-nr}
% ============================================================

\subsubsection{Mathematical Model}
\label{sec:3ph-nr-math}

\paragraph{Network model.}
The three-phase unbalanced network is represented by a compact
$abc$-domain admittance matrix $\mathbf{Y}_{\text{bus}} \in
\mathbb{C}^{n_\phi \times n_\phi}$, where $n_\phi \le 3N$ is the
number of \emph{active} phase-nodes ($N$ buses, each present on up to
three phases~$a, b, c$).  Phase membership is encoded by a bitmask
$\mathcal{P}_i \subseteq \{a,b,c\}$ for each bus~$i$; only the nodes
in $\bigcup_i \mathcal{P}_i$ appear in $\mathbf{Y}_{\text{bus}}$, giving
a reduced problem that handles single-phase laterals and two-phase branches
exactly without phantom nodes.

Each three-phase line contributes a $3\times 3$ (or smaller) block to
$\mathbf{Y}_{\text{bus}}$.  The model supports two impedance representations:
\begin{enumerate}
  \item \textbf{Sequence-derived circulant}: positive-sequence impedance
    $z_1 = r_1 + jx_1$ and zero-sequence impedance $z_0 = r_0 + jx_0$
    give the symmetrical component matrix
    \[
      \mathbf{Z}_{abc} = \mathbf{T}
      \begin{pmatrix} z_1 & & \\ & z_1 & \\ & & z_0 \end{pmatrix}
      \mathbf{T}^{-1},
    \]
    where $\mathbf{T} = \bigl[\mathbf{a}_0,\, \mathbf{a}_1,\, \mathbf{a}_2\bigr]$
    is the symmetrical components transformation matrix with
    $a = e^{j2\pi/3}$.
  \item \textbf{Explicit phase-domain matrix}: when
    \texttt{use\_phase\_matrix = true}, the full $3\times 3$ pu impedance
    matrix $\mathbf{Z}^{abc}$ is used directly, enabling any
    asymmetric mutual coupling.
\end{enumerate}

\paragraph{State vector.}
The NR state is expressed in polar form:
\[
  \mathbf{x} =
  \bigl[\theta_{a,1},\; \theta_{b,1},\; \theta_{c,1},\;\ldots,\;
        \theta_{a,N},\; \theta_{b,N},\; \theta_{c,N},\;
        |V_{a,1}|,\; |V_{b,1}|,\; |V_{c,1}|,\;\ldots\bigr]^{\top},
\]
where $\theta_{\varphi,i}$ is the voltage angle and $|V_{\varphi,i}|$
the magnitude on phase $\varphi$ at bus~$i$.  Slack bus phases are
excluded from $\theta$ (fixed reference); PQ bus phases contribute
to both the $\theta$ and $|V|$ sub-vectors.  PV bus phases are fixed
in magnitude (set on the bus struct) and contribute only to $\theta$.

\paragraph{Mismatch equations.}
For each active phase-node $(\varphi, i)$ the active and reactive power
mismatches are:
\begin{align}
  \Delta P_{\varphi i} &= P^{\text{spec}}_{\varphi i}
    - \sum_{j}\sum_{\psi}
      |V_{\varphi i}|\,|V_{\psi j}|
      \Bigl(
        G_{\varphi\psi,ij}\cos(\theta_{\varphi i}-\theta_{\psi j})
       +B_{\varphi\psi,ij}\sin(\theta_{\varphi i}-\theta_{\psi j})
      \Bigr),
  \label{eq:3ph-p-mismatch}\\[4pt]
  \Delta Q_{\varphi i} &= Q^{\text{spec}}_{\varphi i}
    - \sum_{j}\sum_{\psi}
      |V_{\varphi i}|\,|V_{\psi j}|
      \Bigl(
        G_{\varphi\psi,ij}\sin(\theta_{\varphi i}-\theta_{\psi j})
       -B_{\varphi\psi,ij}\cos(\theta_{\varphi i}-\theta_{\psi j})
      \Bigr),
  \label{eq:3ph-q-mismatch}
\end{align}
where $G_{\varphi\psi,ij}+jB_{\varphi\psi,ij}$ is the
$(\varphi i,\psi j)$ entry of $\mathbf{Y}_{\text{bus}}$, and the
specified injections are
\[
  P^{\text{spec}}_{\varphi i} =
    P^{\text{gen}}_{\varphi i} - P^{\text{load}}_{\varphi i}, \qquad
  Q^{\text{spec}}_{\varphi i} =
    Q^{\text{gen}}_{\varphi i} - Q^{\text{load}}_{\varphi i}.
\]
ZIP loads are linearised around the operating voltage; delta-connected
loads are handled via a Norton-equivalent current injection before
forming the mismatch.

\paragraph{Jacobian.}
The Jacobian $\mathbf{J} \in \mathbb{R}^{n_\text{var} \times n_\text{var}}$
with $n_\text{var} = n_{\theta} + n_{|V|}$ has the standard four sub-blocks:
\[
  \mathbf{J} =
  \begin{pmatrix}
    \dfrac{\partial\boldsymbol{\Delta P}}{\partial\boldsymbol{\theta}} &
    \dfrac{\partial\boldsymbol{\Delta P}}{\partial|\mathbf{V}|} \\[8pt]
    \dfrac{\partial\boldsymbol{\Delta Q}}{\partial\boldsymbol{\theta}} &
    \dfrac{\partial\boldsymbol{\Delta Q}}{\partial|\mathbf{V}|}
  \end{pmatrix}.
\]
The diagonal and off-diagonal entries follow the standard polar-form
expressions (Arrillaga \& Watson, 2001, Ch.\,7), extended to the
multi-phase case by replacing the scalar admittances
$G_{ij}, B_{ij}$ with their inter-phase counterparts
$G_{\varphi\psi,ij}, B_{\varphi\psi,ij}$.
$\mathbf{J}$ is assembled as a sparse matrix and factorised with
Eigen's~\texttt{SparseLU} (or KLU/UMFPACK when available) at each
iteration.

\paragraph{Voltage Unbalance Factor (VUF).}
After convergence the VUF at bus~$i$ is computed via symmetrical
components:
\[
  \mathrm{VUF}_i = \frac{|V_2|_i}{|V_1|_i} \times 100\,\%,
  \qquad
  \begin{pmatrix}V_0\\V_1\\V_2\end{pmatrix}_{\!i}
  = \mathbf{T}^{-1}
  \begin{pmatrix}V_a\\V_b\\V_c\end{pmatrix}_{\!i}.
\]
The result field \texttt{max\_vuf\_percent} reports
$\max_i \mathrm{VUF}_i$.

% ─────────────────────────────────────────────────────────────────────────────
\subsubsection{Algorithm}
\label{sec:3ph-nr-algo}

\begin{algorithm}[H]
\caption{Three-Phase Newton-Raphson Power Flow (\texttt{solve\_three\_phase\_nr})}
\label{alg:3ph-nr}
\begin{algorithmic}[1]
\Require $\mathcal{S}$ (ThreePhaseACSystem), options $(k_{\max}, \varepsilon)$
\Ensure  $\mathbf{r}$ (ThreePhaseDPFResult)

\State Build \textbf{PhaseNodeIndexer}: map each active $(\varphi, i)$ pair to compact index
\State Assemble compact $\mathbf{Y}_{\text{bus}}$ from lines (sequence or phase-domain)
\State Initialise state $\mathbf{x}^{(0)}$ from bus flat-start or per-phase setpoints
\State Build load model: resolve per-phase PQ, ZIP, delta-wye Norton equivalents
\State Compute specified injections $P^{\text{spec}}, Q^{\text{spec}}$ for all active nodes

\For{$k = 0, 1, \ldots, k_{\max}$}
  \State Evaluate mismatch $\boldsymbol{\Delta}(\mathbf{x}^{(k)})$ via
         Eqs.~(\ref{eq:3ph-p-mismatch})--(\ref{eq:3ph-q-mismatch})
  \State $\rho \leftarrow \|\boldsymbol{\Delta}\|_\infty$ \Comment{convergence check}
  \If{$\rho < \varepsilon$}
    \State \textbf{break}
  \EndIf
  \State Assemble sparse Jacobian $\mathbf{J}(\mathbf{x}^{(k)})$
  \State Solve $\mathbf{J}\,\delta\mathbf{x} = \boldsymbol{\Delta}$
         \hfill (SparseLU / KLU)
  \State $\mathbf{x}^{(k+1)} \leftarrow \mathbf{x}^{(k)} + \delta\mathbf{x}$
  \State Clamp $|V| \ge 0.05$\,pu (projection step)
\EndFor

\State Unpack $\mathbf{x}^{(k)} \to |V_{\varphi i}|,\; \theta_{\varphi i}$
\State Compute per-phase branch power flows and losses
\State Compute transformer terminal observations
\State Compute VUF at each bus $\to$ \texttt{max\_vuf\_percent}
\State \Return $\mathbf{r}$ with \texttt{bus\_voltages}, \texttt{branch\_powers},
               \texttt{converged}, \texttt{iterations}, \texttt{residual}
\end{algorithmic}
\end{algorithm}

% ─────────────────────────────────────────────────────────────────────────────
\subsubsection{Flowchart}
\label{sec:3ph-nr-flowchart}

\begin{figure}[H]
\centering
\begin{tikzpicture}[
    node distance = 6mm and 32mm,
    box/.style = {rectangle, draw, rounded corners=3pt,
                  minimum width=52mm, minimum height=8mm,
                  text centered, font=\small},
    dec/.style = {diamond, draw, aspect=2.4,
                  minimum width=52mm, minimum height=8mm,
                  text centered, font=\small},
    arr/.style = {-Latex, thick},
    every node/.append style={align=center}
  ]

  \node[box]             (init)    {Build PhaseNodeIndexer,\\
                                    assemble $\mathbf{Y}_{\text{bus}}$,\\
                                    init flat start $\mathbf{x}^{(0)}$};
  \node[box, below=of init] (load)  {Resolve loads: PQ, ZIP,\\
                                     delta Norton, transformers};
  \node[box, below=of load] (spec)  {Compute $P^{\text{spec}},Q^{\text{spec}}$};
  \node[dec, below=of spec] (iter)  {$k \le k_{\max}$?};
  \node[box, below=of iter] (mism)  {Evaluate mismatch\\
                                     $\boldsymbol{\Delta}(\mathbf{x}^{(k)})$};
  \node[dec, below=of mism] (conv)  {$\|\boldsymbol{\Delta}\|_\infty < \varepsilon$?};
  \node[box, below=of conv] (jac)   {Assemble sparse\\
                                     Jacobian $\mathbf{J}$};
  \node[box, below=of jac]  (solve) {Solve $\mathbf{J}\,\delta\mathbf{x}=\boldsymbol{\Delta}$\\
                                     (SparseLU / KLU)};
  \node[box, below=of solve](upd)   {$\mathbf{x} \leftarrow \mathbf{x}+\delta\mathbf{x}$,\\
                                     clamp $|V|\ge 0.05$};
  \node[box, right=of conv, xshift=4mm] (post)  {Unpack voltages,\\
                                     compute branch\\
                                     power \& losses,\\
                                     VUF};

  \draw[arr] (init) -- (load);
  \draw[arr] (load) -- (spec);
  \draw[arr] (spec) -- (iter);
  \draw[arr] (iter) -- node[right, font=\scriptsize]{yes} (mism);
  \draw[arr] (mism) -- (conv);
  \draw[arr] (conv) -- node[right, font=\scriptsize]{no} (jac);
  \draw[arr] (jac)  -- (solve);
  \draw[arr] (solve)-- (upd);
  % loop back
  \draw[arr] (upd.west) -- ++(-14mm,0) |- (iter.west);
  % converged exit
  \draw[arr] (conv.east) -- node[above, font=\scriptsize]{yes} (post.west);
  % diverged exit
  \draw[arr] (iter.east) -- ++(12mm,0) |- node[above right,
              font=\scriptsize, pos=0.25]{no (diverged)} (post.north);
\end{tikzpicture}
\caption{Three-phase Newton-Raphson power flow algorithm flow. The loop
iterates over the outer \texttt{k} counter; early exit occurs when the
maximum mismatch $\|\boldsymbol{\Delta}\|_\infty$ falls below tolerance
$\varepsilon$.  Both normal convergence and iteration-limit exits feed
the post-processing block that computes branch powers and the VUF.}
\label{fig:3ph-nr-flowchart}
\end{figure}

% ─────────────────────────────────────────────────────────────────────────────
\subsubsection{Solver Options and Result Fields}
\label{sec:3ph-nr-options}

\begin{center}
\small
\begin{tabular}{lll}
\hline\hline
Field & Type & Description \\
\hline
\multicolumn{3}{l}{\textbf{ThreePhaseNROptions}} \\
\texttt{max\_iter}         & int    & Maximum NR iterations (default 50) \\
\texttt{max\_control\_iter}& int    & Maximum outer regulator-control iterations (100) \\
\texttt{tol}               & double & Convergence tolerance on
                                       $\|\boldsymbol{\Delta}\|_\infty$ (1e-6) \\
\texttt{verbose}           & bool   & Print per-iteration residual \\
\texttt{include\_shunts}   & bool   & Add bus shunt $g_s + jb_s$ to $\mathbf{Y}_{\text{bus}}$ \\
\hline
\multicolumn{3}{l}{\textbf{ThreePhaseDPFResult}} \\
\texttt{bus\_voltages}     & vector & Per-bus per-phase $(|V|, \theta)$ \\
\texttt{branch\_powers}    & vector & Per-branch per-phase sending-end $P, Q$ (MW/MVAr) \\
\texttt{p\_loss\_[a/b/c]\_mw}   & double & Per-phase active loss (MW) \\
\texttt{q\_loss\_[a/b/c]\_mvar} & double & Per-phase reactive loss (MVAr) \\
\texttt{total\_p\_loss\_mw}     & double & $\sum_\varphi P^\text{loss}_\varphi$ \\
\texttt{total\_q\_loss\_mvar}   & double & $\sum_\varphi Q^\text{loss}_\varphi$ \\
\texttt{max\_vuf\_percent}      & double & Maximum voltage unbalance factor (\%) \\
\texttt{converged}         & bool   & \texttt{true} if NR converged \\
\texttt{iterations}        & int    & NR iterations taken \\
\texttt{residual}          & double & Final $\|\boldsymbol{\Delta}\|_\infty$ \\
\hline\hline
\end{tabular}
\end{center}

% ─────────────────────────────────────────────────────────────────────────────
\subsubsection{Numerical Verification}
\label{sec:3ph-nr-results}

\paragraph{Balanced 3-bus radial feeder.}
A balanced three-bus radial system ($S_{\text{base}} = 10\,\text{MVA}$,
line $r + jx = 0.04 + j0.08$\,pu) with equal per-phase loads
($P_\varphi = 0.3$\,MW, $Q_\varphi = 0.1$\,MVAr at bus~2;
 $P_\varphi = 0.2$\,MW, $Q_\varphi = 0.08$\,MVAr at bus~3)
converges in $\leq 5$ iterations with $\|\boldsymbol{\Delta}\|_\infty < 10^{-8}$.
All per-phase magnitudes are equal at each bus:
$|V_{a,i}| = |V_{b,i}| = |V_{c,i}|$, consistent with balanced loading.
The VUF is below $0.5\,\%$.

\paragraph{Unbalanced feeder (asymmetric loads).}
With loads concentrated on phase~A (0.5\,MW/0.2\,MVAr at bus~2),
phase~A exhibits the greatest voltage drop.  The NR solver still
converges, and $|V_{a,2}| < |V_{b,2}|$, correctly reflecting the
heavier loading.

\paragraph{Meshed topology.}
A ring branch from bus~3 back to bus~1 ($r + jx = 0.06 + j0.12$\,pu)
creates a meshed network.  The three-phase NR converges in all test cases; all per-phase voltages
remain in the range $(0.85, 1.0)$\,pu.

\paragraph{Single-phase lateral.}
Bus~3 connected via a phase-A-only branch returns
$|V_{b,3}| = |V_{c,3}| = 0$ (no phase-B/C nodes present at that bus),
and all branch-power entries for phases~B and~C are identically zero.

\paragraph{Zero-load (no-load) case.}
With all demand set to zero, all bus voltages converge to exactly
$1.0$\,pu and total losses are zero to machine precision.

\noindent
All 16 test cases in \texttt{tests/test\_three\_phase\_nr.cpp}
(97 assertions) pass with zero failures.
